% concatenated from LHopitalRuleArxivFeb2026.tex, macro.tex
%Updated and Revised Monthly Template 5.7.2024 by Bonnie Ponce and Bev Ruedi copyright the Mathematical Association of America and all rights reserved.

\documentclass{article}

%% For submission to AMM
% \def\AMM{}

\ifdefined\AMM
\usepackage{maa-monthly}
\else
% \usepackage{amssymb}
% \usepackage{amsthm}
% \usepackage{mathrsfs}
\usepackage{maa-monthly-arxiv}
\fi
% \theoremstyle{theorem}
% \newtheorem{theorem}{Theorem}
% \newtheorem{proposition}[theorem]{Proposition}
% \newtheorem{lemma}[theorem]{Lemma}
% \newtheorem{corollary}[theorem]{Corollary}

% \theoremstyle{definition}
% \newtheorem{definition}[theorem]{Definition}
% \newtheorem{remark}[theorem]{Remark}

%============================
% Packages
%============================



%%%%%% Graphing %%%%%%%%%
\usepackage{picinpar,moresize,xfrac,graphpap,dcolumn,wrapfig,graphicx}
\def\figurewtext{1}
%\microtypesetup{stretch=30,shrink=30}
\DeclareGraphicsExtensions{.png,.pdf,.jpg}
\usepackage{subcaption}
\usepackage{stackengine}

%%%%% packages for tikz plots %%%%%
\usepackage{tikz,pgfplots}
\usepackage{tikz-cd}
\usepackage{pgf}
\usepgfplotslibrary{colormaps}
\usetikzlibrary{pgfplots.colormaps}
\usetikzlibrary{arrows,automata,positioning,backgrounds}
% \usetikzlibrary{shapes.geometric,backgrounds,
%   positioning-plus,node-families,calc}
\usepackage{rotating}
\pgfplotsset{compat=newest}
\pgfplotsset{plot coordinates/math parser=false}
\newlength\figureheight
\newlength\figurewidth
%%%%% end of tikz packages %%%%%%%%

%%%%% packages for comments %%%%%
\usepackage[draft]{pdfcomment}
\usepackage{hyperref}
\hypersetup{
    colorlinks=true,
    linkcolor=black,
    filecolor=magenta,      
    urlcolor=magenta,
}

\usepackage[subrefformat=parens]{subcaption}
\captionsetup{compatibility=false}

\definecolor{revcolor}{rgb}{0.0, 0.3, 0.8}

% \newcommand{\rev}[1]{\begingroup\color{revcolor}#1\endgroup}

% \newcommand{\COMMENT}[3]{\textcolor{#2}{{[\textbf{#1:} \textsl{#3}}]}}  % For 

\newcommand{\COMMENT}[3]{\begingroup\color{#2}\textbf{#1: }#3\endgroup}

% work in progress using in-text comments
\newcommand{\TODO}[1]{\COMMENT{TODO: }{red}{#1}}
\newcommand{\Sadashige}[1]{\COMMENT{Sadashige: }{blue}{#1}}
\newcommand{\Albert}[1]{{\color[rgb]{0.2,0.4,0.8}\textbf{Albert: } #1}}

\newcommand{\suggest}[1]{{\color[rgb]{0.0,0.3,0.8} #1}}

\newcommand{\Old}[1]{{\color[rgb]{0.4,0.4,0.4}\textbf{Old: } #1}}



%%%%% End of comment packages %%%%%

%%%%% additional stuff %%%%%%
\usepackage{soul,array,calc,url,ragged2e,graphpap}
 \urlstyle{rm}
\usepackage{booktabs} % better tabular
\usepackage{tabularx}
\usepackage{colortbl}
\usepackage{multirow}
\usepackage{hyphenat}
\usepackage{transparent}
\usepackage{enumerate}
\usepackage{mathrsfs}
% \usepackage{showkeys}
\usepackage{mdframed}
\usepackage{acro}

\usepackage[figure]{hypcap}
%\usepackage{mathabx}
\usepackage{mathtools}
\mathtoolsset{centercolon}
\usepackage{nicefrac}
%\usepackage{siunitx}
\usepackage{units}
% strikethrough machinery
\usepackage[normalem]{ulem}
\usepackage{cancel}
\usepackage{pbox}
\usepackage{multicol}
\usepackage{cases}
\usepackage[all]{xy}
\usepackage{nicematrix}

\usepackage{adjustbox}
\usepackage{wrapfig}
\AfterEndEnvironment{wrapfigure}{\setlength{\intextsep}{0mm}}

\newcommand\AddLabel[1]{\refstepcounter{equation}(\theequation)\label{#1}}
\usepackage{cleveref}
\crefmultiformat{subequation}%
{\edef\crefstripprefixinfo{#1}(#2#1#3}%
{,#2\crefstripprefix{\crefstripprefixinfo}{#1}#3)}%
{,#2\crefstripprefix{\crefstripprefixinfo}{#1}#3}%
{,#2\crefstripprefix{\crefstripprefixinfo}{#1}#3)}

%============================
% Conjugate, Double bracket
%============================
\newcommand{\overbar}[1]{\mkern 1.5mu\overline{\mkern-1.5mu#1\mkern-1.5mu}\mkern 1.5mu}
\newcommand*{\conj}[1]{\overbar{#1}}

\makeatletter
\DeclareFontFamily{OMX}{MnSymbolE}{}
\DeclareSymbolFont{MnLargeSymbols}{OMX}{MnSymbolE}{m}{n}
\SetSymbolFont{MnLargeSymbols}{bold}{OMX}{MnSymbolE}{b}{n}
\DeclareFontShape{OMX}{MnSymbolE}{m}{n}{
    <-6>  MnSymbolE5
   <6-7>  MnSymbolE6
   <7-8>  MnSymbolE7
   <8-9>  MnSymbolE8
   <9-10> MnSymbolE9
  <10-12> MnSymbolE10
  <12->   MnSymbolE12
}{}
\DeclareFontShape{OMX}{MnSymbolE}{b}{n}{
    <-6>  MnSymbolE-Bold5
   <6-7>  MnSymbolE-Bold6
   <7-8>  MnSymbolE-Bold7
   <8-9>  MnSymbolE-Bold8
   <9-10> MnSymbolE-Bold9
  <10-12> MnSymbolE-Bold10
  <12->   MnSymbolE-Bold12
}{}
\let\llangle\@undefined
\let\rrangle\@undefined
\DeclareMathDelimiter{\llangle}{\mathopen}%
                     {MnLargeSymbols}{'164}{MnLargeSymbols}{'164}
\DeclareMathDelimiter{\rrangle}{\mathclose}%
                     {MnLargeSymbols}{'171}{MnLargeSymbols}{'171}
\makeatother

%==========================================
% horizontal and vertical lines in matrices
%==========================================
\newcommand*{\vertbar}{\rule[-1ex]{0.5pt}{2.5ex}}
\newcommand*{\horzbar}{\rule[.5ex]{2.5ex}{0.5pt}}


%======================
% Code and Algorithms
%======================

\usepackage{fancyvrb,listings}
\usepackage{algorithm}
\usepackage{algorithmicx}
\usepackage[noend]{algpseudocode}
\renewcommand{\algorithmicrequire}{\textbf{Input:}}
\renewcommand{\algorithmicensure}{\textbf{Output:}}

\newcommand{\Proc}[1]{\mbox{\textsc{#1}}}
\renewcommand{\algorithmicforall}{\textbf{for each}}
\renewcommand{\algorithmicrepeat}{\textbf{do}}
\renewcommand{\algorithmicuntil}{\textbf{while}}
\let\ForEach\ForAll
\let\Do\Repeat
\let\dWhile\Until

\algrenewcommand\alglinenumber[1]{\footnotesize #1:}
\makeatletter  
 \renewcommand{\ALG@name}{\small Algorithm} 
\makeatother 
\algdef{SE}[DOWHILE]{Do}{doWhile}{\algorithmicdo}[1]{\algorithmicwhile\ #1}%

%\renewcommand{\thealgorithm}{\arabic{section}.\arabic{algorithm}}

%================
% Theorems
%================
\theoremstyle{theorem}
\newtheorem{theorem}{Theorem}
\newtheorem{lemma}{Lemma}
\newtheorem{proposition}{Proposition}
\newtheorem{corollary}{Corollary}
\newtheorem{problem}{Problem}
\newtheorem{example}{Example}
\theoremstyle{definition}
\newtheorem{definition}{Definition}
\newtheorem{question}{Open problem}

\newenvironment{claim}
{\par\medskip
 \phantomsection
 \noindent\textit{Claim.}\itshape
}
{\par\medskip}

% \theoremstyle{remark}
\newtheorem{remark}{Remark}



%================
% REFERENCING
%================
\newcommand{\figref}[1]{\textup{Fig.~\ref{#1}}}
\newcommand{\tabref}[1]{\textup{Table~\ref{#1}}}
\newcommand{\teqref}[1]{\textup{Eq.~(\ref{#1})}}
\newcommand{\chref}[1]{\textup{Chapter~\ref{#1}}}
\newcommand{\secref}[1]{\textup{Section~\ref{#1}}}
\newcommand{\appref}[1]{\textup{Appendix~\ref{#1}}}
\renewcommand{\algref}[1]{\textup{Alg.~\ref{#1}}}
\newcommand{\thmref}[1]{\textup{Theorem~\ref{#1}}}
\newcommand{\lemref}[1]{\textup{Lemma~\ref{#1}}}
\newcommand{\remref}[1]{\textup{Remark~\ref{#1}}}
\newcommand{\defref}[1]{\textup{Definition~\ref{#1}}}
\newcommand{\propref}[1]{\textup{Proposition~\ref{#1}}}
\newcommand{\corref}[1]{\textup{Corollary~\ref{#1}}}
\newcommand{\exref}[1]{\textup{Example~\ref{#1}}}
\newcommand{\exerref}[1]{\textup{Exercise~\ref{#1}}}
\newcommand{\probref}[1]{\textup{Problem~\ref{#1}}}
\newcommand{\notref}[1]{\textup{Notation~\ref{#1}}}




%=================== 
% LATINS
%===================
\def\etal{\emph{et al.}}
\def\cf{\emph{cf.}}
\def\ie{\emph{i.e.}}
\def\eg{\emph{e.g.}}
\def\etc{\emph{etc.}}
\def\etcnp{\emph{etc}}

%===================
% FONT SHORTHANDS
%===================
\def\AA{\mathbb{A}}
\def\BB{\mathbb{B}}
\def\CC{\mathbb{C}}
\def\DD{\mathbb{D}}
\def\EE{\mathbb{E}}
\def\FF{\mathbb{F}}
\def\GG{\mathbb{G}}
\def\HH{\mathbb{H}}
\def\II{\mathbb{I}}
\def\JJ{\mathbb{J}}
\def\KK{\mathbb{K}}
\def\LL{\mathbb{L}}
\def\MM{\mathbb{M}}
\def\NN{\mathbb{N}}
\def\OO{\mathbb{O}}
\def\PP{\mathbb{P}}
\def\QQ{\mathbb{Q}}
\def\RR{\mathbb{R}}
\def\SS{\mathbb{S}}
\def\TT{\mathbb{T}}
\def\UU{\mathbb{U}}
\def\VV{\mathbb{V}}
\def\WW{\mathbb{W}}
\def\XX{\mathbb{X}}
\def\YY{\mathbb{Y}}
\def\ZZ{\mathbb{Z}}

\def\bA{\mathbf{A}}
\def\bB{\mathbf{B}}
\def\bC{\mathbf{C}}
\def\bD{\mathbf{D}}
\def\bE{\mathbf{E}}
\def\bF{\mathbf{F}}
\def\bG{\mathbf{G}}
\def\bH{\mathbf{H}}
\def\bI{\mathbf{I}}
\def\bJ{\mathbf{J}}
\def\bK{\mathbf{K}}
\def\bL{\mathbf{L}}
\def\bM{\mathbf{M}}
\def\bN{\mathbf{N}}
\def\bO{\mathbf{O}}
\def\bP{\mathbf{P}}
\def\bQ{\mathbf{Q}}
\def\bR{\mathbf{R}}
\def\bS{\mathbf{S}}
\def\bT{\mathbf{T}}
\def\bU{\mathbf{U}}
\def\bV{\mathbf{V}}
\def\bW{\mathbf{W}}
\def\bX{\mathbf{X}}
\def\bY{\mathbf{Y}}
\def\bZ{\mathbf{Z}}

\def\cA{\mathcal{A}}
\def\cB{\mathcal{B}}
\def\cC{\mathcal{C}}
\def\cD{\mathcal{D}}
\def\cE{\mathcal{E}}
\def\cF{\mathcal{F}}
\def\cG{\mathcal{G}}
\def\cH{\mathcal{H}}
\def\cI{\mathcal{I}}
\def\cJ{\mathcal{J}}
\def\cK{\mathcal{K}}
\def\cL{\mathcal{L}}
\def\cM{\mathcal{M}}
\def\cN{\mathcal{N}}
\def\cO{\mathcal{O}}
\def\cP{\mathcal{P}}
\def\cQ{\mathcal{Q}}
\def\cR{\mathcal{R}}
\def\cS{\mathcal{S}}
\def\cT{\mathcal{T}}
\def\cU{\mathcal{U}}
\def\cV{\mathcal{V}}
\def\cW{\mathcal{W}}
\def\cX{\mathcal{X}}
\def\cY{\mathcal{Y}}
\def\cZ{\mathcal{Z}}

\def\sA{\mathscr{A}}
\def\sB{\mathscr{B}}
\def\sC{\mathscr{C}}
\def\sD{\mathscr{D}}
\def\sE{\mathscr{E}}
\def\sF{\mathscr{F}}
\def\sG{\mathscr{G}}
\def\sH{\mathscr{H}}
\def\sI{\mathscr{I}}
\def\sJ{\mathscr{J}}
\def\sK{\mathscr{K}}
\def\sL{\mathscr{L}}
\def\sM{\mathscr{M}}
\def\sN{\mathscr{N}}
\def\sO{\mathscr{O}}
\def\sP{\mathscr{P}}
\def\sQ{\mathscr{Q}}
\def\sR{\mathscr{R}}
\def\sS{\mathscr{S}}
\def\sT{\mathscr{T}}
\def\sU{\mathscr{U}}
\def\sV{\mathscr{V}}
\def\sW{\mathscr{W}}
\def\sX{\mathscr{X}}
\def\sY{\mathscr{Y}}
\def\sZ{\mathscr{Z}}

\def\sfA{\mathsf{A}}
\def\sfB{\mathsf{B}}
\def\sfC{\mathsf{C}}
\def\sfD{\mathsf{D}}
\def\sfE{\mathsf{E}}
\def\sfF{\mathsf{F}}
\def\sfG{\mathsf{G}}
\def\sfH{\mathsf{H}}
\def\sfI{\mathsf{I}}
\def\sfJ{\mathsf{J}}
\def\sfK{\mathsf{K}}
\def\sfL{\mathsf{L}}
\def\sfM{\mathsf{M}}
\def\sfN{\mathsf{N}}
\def\sfO{\mathsf{O}}
\def\sfP{\mathsf{P}}
\def\sfQ{\mathsf{Q}}
\def\sfR{\mathsf{R}}
\def\sfS{\mathsf{S}}
\def\sfT{\mathsf{T}}
\def\sfU{\mathsf{U}}
\def\sfV{\mathsf{V}}
\def\sfW{\mathsf{W}}
\def\sfX{\mathsf{X}}
\def\sfY{\mathsf{Y}}
\def\sfZ{\mathsf{Z}}

\def\sfa{\mathsf{a}}
\def\sfb{\mathsf{b}}
\def\sfc{\mathsf{c}}
\def\sfd{\mathsf{d}}
\def\sfe{\mathsf{e}}
\def\sff{\mathsf{f}}
\def\sfg{\mathsf{g}}
\def\sfh{\mathsf{h}}
\def\sfi{\mathsf{i}}
\def\sfj{\mathsf{j}}
\def\sfk{\mathsf{k}}
\def\sfl{\mathsf{l}}
\def\sfm{\mathsf{m}}
\def\sfn{\mathsf{n}}
\def\sfo{\mathsf{o}}
\def\sfp{\mathsf{p}}
\def\sfq{\mathsf{q}}
\def\sfr{\mathsf{r}}
\def\sfs{\mathsf{s}}
\def\sft{\mathsf{t}}
\def\sfu{\mathsf{u}}
\def\sfv{\mathsf{v}}
\def\sfw{\mathsf{w}}
\def\sfx{\mathsf{x}}
\def\sfy{\mathsf{y}}
\def\sfz{\mathsf{z}}


\def\ba{\mathbf{a}}
\def\bb{\mathbf{b}}
\def\bc{\mathbf{c}}
\def\bd{\mathbf{d}}
\def\be{\mathbf{e}}
\def\bff{\mathbf{f}}
\def\bg{\mathbf{g}}
\def\bh{\mathbf{h}}
\def\bi{\mathbf{i}}
\def\bj{\mathbf{j}}
\def\bk{\mathbf{k}}
\def\bl{\mathbf{l}}
\def\bm{\mathbf{m}}
\def\bn{\mathbf{n}}
\def\bo{\mathbf{o}}
\def\bp{\mathbf{p}}
\def\bq{\mathbf{q}}
\def\br{\mathbf{r}}
\def\bs{\mathbf{s}}
\def\bt{\mathbf{t}}
\def\bu{\mathbf{u}}
\def\bv{\mathbf{v}}
\def\bw{\mathbf{w}}
\def\bx{\mathbf{x}}
\def\by{\mathbf{y}}
\def\bz{\mathbf{z}}

\def\balpha{\boldsymbol{\alpha}}
\def\bbeta{\boldsymbol{\beta}}
\def\bgamma{\boldsymbol{\gamma}}
\def\bGamma{\boldsymbol{\Gamma}}
\def\bdelta{\boldsymbol{\delta}}
\def\bDelta{\boldsymbol{\Delta}}
\def\bepsilon{\boldsymbol{\epsilon}}
\def\bvarepsilon{\boldsymbol{\varepsilon}}
\def\bzeta{\boldsymbol{\zeta}}
\def\boldeta{\boldsymbol{\eta}}
\def\btheta{\boldsymbol{\theta}}
\def\bvartheta{\boldsymbol{\vartheta}}
\def\bTheta{\boldsymbol{\Theta}}
\def\biota{\boldsymbol{\iota}}
\def\bkappa{\boldsymbol{\kappa}}
\def\blambda{\boldsymbol{\lambda}}
\def\bLambda{\boldsymbol{\Lambda}}
\def\bmu{\boldsymbol{\mu}}
\def\bnu{\boldsymbol{\nu}}
\def\bxi{\boldsymbol{\xi}}
\def\bXi{\boldsymbol{\Xi}}
\def\bpi{\boldsymbol{\pi}}
\def\bvarpi{\boldsymbol{\varpi}}
\def\bPi{\boldsymbol{\Pi}}
\def\brho{\boldsymbol{\rho}}
\def\bvarrho{\boldsymbol{\varrho}}
\def\bsigma{\boldsymbol{\sigma}}
\def\bvarsigma{\boldsymbol{\varsigma}}
\def\bSigma{\boldsymbol{\Sigma}}
\def\btau{\boldsymbol{\tau}}
\def\bphi{\boldsymbol{\phi}}
\def\bvarphi{\boldsymbol{\varphi}}
\def\bPhi{\boldsymbol{\Phi}}
\def\bchi{\boldsymbol{\chi}}
\def\bpsi{\boldsymbol{\psi}}
\def\bPsi{\boldsymbol{\Psi}}
\def\bomega{\boldsymbol{\omega}}
\def\bOmega{\boldsymbol{\Omega}}

\def\fB{\mathfrak{B}}
\def\fD{\mathfrak{D}}
\def\fH{\mathfrak{H}}
\def\fX{\mathfrak{X}}
\def\fg{\mathfrak{g}}

\def\bzero{\mathbf{0}}
\def\bpartial{\boldsymbol{\partial}}

\def\continuity{\mathcal{C}}
\def\dt{{\Deltait t}}
\def\adjoint{\intercal}

%=====================
% QUATERNIONS
%=====================
\usepackage{bbm}
%%%%%makes a stylized 1%%%%%%%%%
\DeclareSymbolFont{bbold}{U}{bbold}{m}{n}
\DeclareSymbolFontAlphabet{\mathbbold}{bbold}
\newcommand{\one}{\mathbbold{1}}%%%%%%%%%makes stylized 1
% \newcommand{\ii}{\mathbbm i}
% \newcommand{\jj}{\mathbbm j}
% \newcommand{\kk}{\mathbbm k}
\newcommand{\ii}{\mkern1.5mu\mathbbold{i}\mkern1.5mu}
\newcommand{\jj}{\mkern1.5mu\mathbbold{j}\mkern1.5mu}
\newcommand{\kk}{\mkern1.5mu\mathbbold{k}\mkern1.5mu} 
%%%%%%%%%%%%%%%%%%%%

%=====================
% PROJECTIVE SPACES
%=====================
\newcommand{\RP}{\mathbb{RP}}
\newcommand{\CP}{\mathbb{CP}}
\newcommand{\HP}{\mathbb{HP}}

%=====================
% MATH OPERATORS
%=====================

% real, imaginary, and other parts of numbers
\renewcommand{\Im}{\operatorname{Im}}
\renewcommand{\Re}{\operatorname{Re}}
\newcommand{\sgn}{\operatorname{sgn}}
\DeclareMathOperator{\atan}{atan2}

% linear algebra
\def\tr {\operatorname{tr}}
\def\det{\operatorname{det}}
\def\End{\operatorname{End}}
\def\Hom{\operatorname{Hom}}
\DeclareMathOperator{\Map}{Map}
\def\im {\operatorname{im}}
\def\ker{\operatorname{ker}}
\def\spanset{\operatorname{span}}
\def\id{\operatorname{id}}
\def\adj{\operatorname{adj}}
\def\cof{\operatorname{cof}}
\def\J{\cJ}
\def\botoplus{\mathbin{\overset{\bot}\oplus}}
\def\CrossRatio{\operatorname{cr}}


% optimization
\DeclareMathOperator*{\argmin}{argmin}
\DeclareMathOperator*{\argmax}{argmax}
\DeclareMathOperator*{\arginf}{arginf}
\DeclareMathOperator*{\argsup}{argsup}

% Lie groups and Lie algebras
\DeclareMathOperator{\GL}{GL}
\DeclareMathOperator{\SL}{SL}
\DeclareMathOperator{\PSL}{PSL}
\DeclareMathOperator{\PSU}{PSU}
\DeclareMathOperator{\SO}{SO}
\DeclareMathOperator{\SU}{SU}
\renewcommand{\so}{\mathfrak{so}} % \so already defined.
\DeclareMathOperator{\su}{\mathfrak{su}}
\DeclareMathOperator{\gl}{\mathfrak{gl}}
\DeclareMathOperator{\slla}{\mathfrak{sl}}
\DeclareMathOperator{\Diff}{Diff}
\DeclareMathOperator{\SDiff}{SDiff}
\DeclareMathOperator{\Spin}{Spin}
\DeclareMathOperator{\Iso}{Iso}
\DeclareMathOperator{\Isom}{Isom}
\DeclareMathOperator{\Aut}{Aut}
\DeclareMathOperator{\Conj}{Conj}
\DeclareMathOperator{\Ad}{Ad}
\DeclareMathOperator{\ad}{ad}
\DeclareMathOperator{\diff}{\mathfrak{diff}}
\DeclareMathOperator{\sdiff}{\mathfrak{sdiff}}
\DeclareMathOperator{\isom}{\mathfrak{isom}}
\DeclareMathOperator{\Ham}{Ham}
\DeclareMathOperator{\ham}{\mathfrak{ham}}
\DeclareMathOperator{\Symp}{Symp}
\DeclareMathOperator{\symp}{\mathfrak{symp}}

\DeclareMathOperator{\Flux}{\textsc{Flux}}
\DeclareMathOperator{\Transf}{\textsc{Transf}}
\DeclareMathOperator{\Link}{\textsc{Link}}
\DeclareMathOperator{\Div}{Div}
\DeclareMathOperator{\ord}{ord}
\DeclareMathOperator{\Cl}{Cl}
\DeclareMathOperator{\Pic}{Pic}
\DeclareMathOperator{\AJ}{AJ}

\DeclareMathOperator{\Flow}{Fl}
\DeclareMathOperator{\Relab}{Rl}
\DeclareMathOperator{\flow}{fl}
\DeclareMathOperator{\relab}{rl}
\DeclareMathOperator{\Der}{Der}

\DeclareMathOperator{\Act}{Act}
\DeclareMathOperator{\act}{act}

\DeclareMathOperator{\Adv}{Adv}
\DeclareMathOperator{\adv}{adv}

\DeclarePairedDelimiter\parens{\lparen}{\rparen}
\DeclarePairedDelimiter\abs{\lvert}{\rvert}
\DeclarePairedDelimiter\norm{\lVert}{\rVert}
\DeclarePairedDelimiter\floor{\lfloor}{\rfloor}
\DeclarePairedDelimiter\ceil{\lceil}{\rceil}
\DeclarePairedDelimiter\braces{\lbrace}{\rbrace}
\DeclarePairedDelimiter\bracks{\lbrack}{\rbrack}
\DeclarePairedDelimiter\angles{\langle}{\rangle}

\DeclarePairedDelimiter\bra{\langle}{\rvert}
\DeclarePairedDelimiter\ket{\lvert}{\rangle}
\DeclarePairedDelimiterX\braket[2]{\langle}{\rangle}{#1\,\delimsize\vert\,\mathopen{}#2}

% Second fundamental form
\DeclareMathOperator{\SFF}{I\!I}

% Vector Calculus
\newcommand{\grad}{\mathop{\mathrm{grad}}\nolimits}
\newcommand{\sgrad}{\mathop{\mathrm{sgrad}}\nolimits}
\newcommand{\curl}{\mathop{\mathrm{curl}}\nolimits}
\newcommand{\rot}{\mathop{\mathrm{rot}}\nolimits}
\renewcommand{\div}{\mathop{\mathrm{div}}\nolimits}
\newcommand{\LD}{\mathop{\mathscr{L}}\nolimits}
\newcommand{\Hess}{\mathop{\mathrm{Hess}}\nolimits}

\newcommand{\checkgrad}{\mathop{\check{\smash{\mathrm{grad}}}}\nolimits}
\newcommand{\checkdiv}{\mathop{\check{{\mathrm{div}}}}\nolimits}

% Probability
\def\Cov{\operatorname{Cov}}
\def\Var{\operatorname{Var}}

% Exterior Calculus
\def\wedgetimes{\mathbin{\underset{\times}{\wedge}}}
\def\wedgecomma{\mathbin{\hspace{0.5ex}\raisebox{-0.5ex}{,}\hspace{-1ex}{\wedge}}}

% Interior Calculus
\def\tt{\mathbbold{t}}
\def\nn{\mathbbold{n}}
\def\contract{\lrcorner}
\def\ip{i}
\def\ext{\operatorname{ext}}

% Geometry
\DeclareMathOperator{\Area}{Area}
\DeclareMathOperator{\Length}{Length}
\DeclareMathOperator{\Vol}{Vol}
\DeclareMathOperator{\Dihedral}{Dihedral}
\def\cp{\mathbbold{cp}}

% Numerics
\def\timestep{{\Deltait} t}

% Convention
% Sets
\def\invOmega{\Omega_{\rm inv}}
\def\invSigma{\Sigma_{\rm inv}}
% Coordinates
\def\invbx{\by}
\def\invx{y}
\def\invy{Y}
\def\invz{Z}
% Functions
\def\invf{F}
\def\invg{G}
\def\invu{U}
\def\invv{V}
\def\invp{P}
\def\invq{Q}
% Vectors
\def\invba{\bA}
\def\invbb{\bB}
\def\invbw{\bW}
% Normal
\def\invbn{\bN}
% Differential
\def\invnabla{\nabla}
\def\invDelta{\Delta}

% Additional bold symbols, operators
\def\bpsi{\boldsymbol{\psi}}
\def\Flux{\textnormal{\textsc{Flux}}}
\def\Link{\textnormal{\textsc{Link}}}
\def\bxi{\boldsymbol{\xi}}
\def\bzeta{\boldsymbol{\zeta}}



\definecolor{b1}{rgb}{0.158099,0.313781,0.636957}
\definecolor{b2}{rgb}{0.525367,0.691857,0.998936}
\definecolor{g1}{rgb}{0.256000,0.640000,0.576000}
\definecolor{g2}{rgb}{0.559573,0.800781,0.760580}
\definecolor{p1}{rgb}{0.416000,0.192000,0.640000}
\definecolor{p2}{rgb}{0.680880,0.561880,0.799881}
\definecolor{r1}{rgb}{0.800000,0.200000,0.000000}
\definecolor{r2}{rgb}{1.000000,0.702595,0.603461}
\definecolor{k1}{gray}{0}
\definecolor{k2}{gray}{0.7}
\newcommand\Cb[1]{\color{b#1}}
\newcommand\Cg[1]{\color{g#1}}
\newcommand\Cp[1]{\color{p#1}}
\newcommand\Cr[1]{\color{r#1}}
\newcommand\Ck[1]{\color{k#1}}
\usepackage{enumitem}

\begin{document}
%Instructions: A title page has been added to this template on 6.18.2025 for the purpose of matching meta data entered into the Editorial Manager System for The American Mathematical Monthly and the manuscript.  Please fill out this title form for both the manuscript with author details and the manuscript without author details. Compile one PDF that has the author details and then comment out author details to compile another PDF that is anonymous or without the author details.



\ifdefined\AMM
\begin{center}
{\LARGE{Title Page for \textit{The American Mathematical Monthly}}}
\end{center}

\vspace{150px}
\begin{center}

\textbf{l'H\^opital rules for complex-valued functions in higher dimensions} \\
\vspace{25px}

%Author 1 Name, University or Affiliation, City, Country, email address; this information should not appear in the manuscript without author details
Albert Chern,\\ 
UCSD,
California, USA, 
alchern@ucsd.edu

\vspace{15px}

%Author 2 Name, University or Affiliation, City, Country, email address; this information should not appear in the manuscript without author details
Sadashige Ishida (corresponding author), \\
ISTA,
Klosterneuburg, Austria, 
sadashige.ishida@ist.ac.at

\vspace{15px}


%copy and past if more authors than 3.
\end{center}
\pagebreak
\fi

\title{l'H\^opital rules for complex-valued functions in higher dimensions} %Put Manuscritp Title Here
\markright{Complex L’H\^opital Rules in Higher Dimensions}
% \author{Double Blind No Name1 and Double Blind No Name2}  %DO NOT FILL IN AUTHOR'S NAMES UNTIL YOU RECEIVE YOUR PROVISIONAL ACCEPT LETTER. SUBMISSIONS TO THE MONTHLY ARE DOUBLE BLIND.
\author{Albert Chern and Sadashige Ishida} 
\maketitle


\begin{abstract}
% We study l'H\^opital-type rules for complex-valued functions in multiple dimensions. We provide a necessary and sufficient condition for the existence of a \emph{continuous} quotient, encoded in an algebraic relation between the derivatives of the input functions.  This condition is surprisingly nontrivial, unlike the case of real-valued functions in multiple dimensions or complex-valued functions in real one dimension. The question of a sharp condition for the existence of a \emph{smooth} quotient remains open.
In calculus, l'H\^opital's rule provides a simple way to evaluate the limits of quotient functions when both the numerator and denominator vanish.  But what happens when we move beyond real functions on a real interval?  In this article, we study when the quotient of two complex-valued functions in higher dimension can be defined continuously at the points where both functions vanish.  Surprisingly, the answer is far subtler than in the real-valued setting.  We provide a complete characterization for the continuity of the quotient function.  We also point out why extending this result to smoother quotients remains an intriguing challenge.

% \Albert{[The current abstract might be too researcher-oriented instead of a tone of ``curiosity and accessibility.''] In calculus, l'H\^opital rule provides a simple way to evaluate the limits of ratios when both the numerator and denominator vanishes.  But what happens when we move beyond real functions of one variable?  In this article, we study when the ratio of two complex-valued functions of several variables can be defined continuously at the points where both functions vanish.  Surprisingly, the answer is far subtler than in the real-valued setting.  We provide a complete characterization for the continuity of the quotient function.  We also point out why extending these results to smoother quotients remains an intriguing challenge.}
% \Sadashige{It's a very nice abstract, thanks!}
\end{abstract}

\vspace{10px}

\noindent{\textbf{Keywords: }{l'H\^opital theorem, complex functions}} 

\vspace{10px}

\section{Introduction}
The classical Bernoulli--l'H\^opital rule for real functions \(f(x)\) and \(g(x)\) on a real interval states conditions under which the ratio \(f(x)/g(x)\) is well-defined at a common zero of \(f\) and \(g\).
Specifically, if \(f(a) = g(a) = 0\), \(g'(a)\neq 0\), and the limit \(L = \lim_{x\to a}f'(x)/g'(x)\) exists, then \(\lim_{x\to a}f(x)/g(x) = L\).

The result was further extended to complex-valued functions on a real interval by \cite{Carter:1958:HRC} and, more recently, to real-valued multivariable functions by \cite{Lawlor:2020:HRM}. 
% Other l'H\^opital style theorems are surveyed in \cite{Lawlor:2020:HRM}.

While the classical l'H\^opital rule and its extensions focus on the limit of \(f(x)/g(x)\) at the common zeros of \(f\) and \(g\), a stronger theorem can be stated in terms of regularity.  If we define \(f(a)/g(a)\coloneqq L\), then \(f(x)/g(x)\) remains a smooth function as long as both \(f\) and \(g\) are smooth. 
\begin{theorem}[Smooth quotient of real functions]
\label{thm:RealQuotient1D}
    Let \(f,g\in C^\infty(I)\) be smooth real-valued functions on a real interval \(I\).  Suppose \(f,g\) have only simple zeros in \(I\) and that their zero sets are identical.  Then there exists a non-vanishing smooth function \(\varphi\in C^\infty(I)\) such that \(f = g\varphi\) and $\varphi=f'/g'$ at the zero level sets.
\end{theorem}
Here, a real function \(f\) on \(I\) is said to have simple zeros if \(f'(a)\neq 0 \) whenever \(f(a) = 0\).
An elementary proof of \autoref{thm:RealQuotient1D} is provided in \autoref{app:FirstOrderTaylorExpansion} and \ref{app:ProofOfRealQuotient1D}.
% \Albert{[This jumps to Appendix B, but then that proof heavily uses Lemma from Appendix A.  It's better to say ``a proof of \autoref{thm:RealQuotient1D} is provided in \autoref{app:ProofOfRealQuotient1D}, which is based a few elementary results about Taylor expansion detailed in \autoref{app:FirstOrderTaylorExpansion}.'']}

Intuitively, functions with common zeros can be thought of as sharing common factors, which cancel out in their quotient.  A similar result holds for real-valued multivariable functions.  Let \(\Omega\subseteq \RR^n\) be an open set, with \(n\geq 2\), serving as the domain.  The zero set of a real-valued function \(f\colon\Omega\to\RR\) is denoted by \(\Sigma_f = \{\ba\in\Omega\,|\, f(\ba) = 0\}\subset\Omega\).
A smooth function \(f\colon\Omega\to\RR\) is said to have simple zeros if \(\nabla f|_\ba \neq 0\) for all \(\ba\in\Sigma_{f}\).  Geometrically, by the implicit function theorem, the zero set \(\Sigma_f\) of a function with simple zeros forms a smooth hypersurface.
\begin{theorem}[Smooth quotient of real functions in arbitrary dimension]
\label{thm:RealQuotientnD}
    Let \(f,g\in C^\infty(\Omega)\) be smooth real-valued functions with simple zeros and a common zero set \(\Sigma_{f} = \Sigma_{g}\).  Then there exists a non-vanishing smooth function \(\varphi\in C^\infty(\Omega)\) such that \(f = g\varphi\).
\end{theorem}
A proof of \autoref{thm:RealQuotientnD} is provided in \autoref{app:ProofOfRealQuotientnD}.

One might ask whether a similar result holds for complex-valued  functions in arbitrary dimension. Specifically, do two complex-valued functions on an open set \(\Omega\subset\RR^n\) with a common zero set always have a smooth complex-valued quotient? The answer to this question is surprisingly different from the real-valued case.

As we will see in Section \ref{sec:RatioOfComplexFunctions}, having common simple zeros is generally insufficient. It turns out that the existence of a continuous quotient requires an additional algebraic relation between the gradients of the functions at the zero set. We describe this condition in the following.

% \TODO{Write something like if $f,g$ are holomorphic around zeros, the quotient trivially exists. But it's a very restrictive condition. Can we say nothing for smooth functions? We want to advance the understanding. 

% In addition, holomorphicity is not invariant under coordinate changes. Can we really say nothing if both $f$ and $g$ become non-holomorphic from being holomorphic under the same coordinate change? Also holomorphicity makes sense only in $2D$. We propose a relaxation, which we call \emph{relative holomorphicity}, which address these issues.}

\section{Quotient of Complex Functions}
\label{sec:RatioOfComplexFunctions}

Most existing generalizations of the l'H\^opital theorem for complex-valued functions are formulated on a real interval \cite{Carter:1958:HRC} or for holomorphic functions on a complex plane, a restrictive class of functions studied in complex analysis \cite{AhlforsLars:1979:complex}.
In the latter case, the quotient $f/g$ is globally defined and it coincides with $f'(\ba)/g'(\ba)$ at each common simple zero $\ba$. 
% \Albert{[This might sound silly, but it's possible that some readers are not familiar with ``holomorphic functions'' if we are targeting to calculus students.  It's also not our job to teach them complex analysis.  I would say ``\emph{holomorphic functions} on complex planes, a rather restrictive class of functions studied in complex analysis'' so that the students would know they will learn it, and the only thing one needs to know now is that they are not general in our context.]}
% If both \(f\) and \(g\) were holomorphic functions with the same simple zeros, then \(f/g\) would be holomorphic and therefore smooth.

% But this condition is very restrictive and somewhat artificial.
Of course, a theory that guarantees the existence of a quotient should be much more general than simply requiring both functions to be locally holomorphic. For instance, the quotient \(f/f = 1\) is smooth even when \(f\) is not holomorphic. 
% In fact, many pairs of non-holomorphic functions admit smooth quotients. For instance, multiplying a nowhere-vanishing smooth function by a (not necessarily holomorphic) smooth function naturally yields a pair $(f,g)$ that admits a smooth quotient. 
Moreover, holomorphy is not invariant under smooth coordinate transformations, whereas the existence of a smooth quotient is independent. Additionally, it is unclear how the notion of holomorphy should make sense in a domain of general real dimension. 
% The notion of holomorphy also does not make sense in real odd dimensions, while smooth quotients are well-defined in any dimension.

So, can we generalize the l'H\^opital theorem for complex-valued functions in a space of arbitrary dimension?  This is the main question we explore in this article. 
% Let us pose this question and discuss its nature in the following.

% So Can we not have any theory for the existence of quotient for non-holormorphic functions in a space of arbitrary dimension? 
% In the present and the next section, we introduce l'H\^opital theorems for such a wider setting. \Albert{[This paragraph reads a bit awkwardly mainly because of the negative tone of ``Can we not''.  Maybe: ``So, can we generalize the l'H\^opital theorem for complex-valued functions in a space of arbitrary dimension?  Let us pose this question and discuss its nature in the following.'']}


Let \(\Omega\subset\RR^n\) be an open set with \(n\geq 2\). For a complex-valued function \(f\colon\Omega\to\CC\), the zero set is defined as \(\Gamma_f \coloneqq \{\ba\in\Omega\,|\, f(\ba) = 0\}\subset\Omega\), 
and the notion of simple zeros for complex-valued functions is characterized by a non-degeneracy condition on the first derivative.

\begin{definition}[Simple zero]
For $n\geq 2$, a differentiable function \(f\colon\Omega(\subset \RR^n) \to\CC\) is said to have \emph{simple zeros} if, at each zero \(\ba\in\Gamma_f\), the derivative \(Df|_{\ba}\colon\RR^n\xrightarrow{\rm linear}\CC\cong\RR^2\), considered as a real linear map, has rank 2. 
\end{definition}

By the implicit function theorem, the zero set \(\Gamma_f\) of a function with simple zeros forms a codimension-2 submanifold embedded in \(\Omega\).

\begin{problem}[Smooth quotient of complex multivariable functions]\label{prob:ComplexLHopital}
Let \(f,g\in C^\infty(\Omega;\CC)\) be smooth complex-valued functions with a common zero set \(\Gamma_f = \Gamma_g\) consisting of simple zeros.  Does there exist a nonvanishing complex-valued function \(\varphi\in C^\infty(\Omega;\CC)\) so that \(f = g\varphi\)?
\end{problem}
In most cases, the answer to \Cref{prob:ComplexLHopital} is negative, as demonstrated in the following examples.
\begin{example}\label{eg:NonexampleRatio1}
    Let \(\Omega=\RR^2\), \(f(x,y)\coloneqq x+y+\ii y\), and \(g(x,y) = x + \ii y\).  They have a common simple zero at the origin.  The quotient of these two functions is given by
    \(
    \varphi(x,y) = \frac{f(x,y)}{g(x,y)} = \frac{x+y+\ii y}{x+\ii y}. 
    \)
    This function is constant along any straight line passing through the origin, yet is undefined at the origin. To see this, let us take a line $(x(t),y(t))=(t,ct)$ with some real constant $c$ and a parameter $t$. Then we have \(\varphi(t,ct)=\frac{t+ct+\ii ct}{t+\ii ct}=\frac{1+c+\ii c}{1+\ii c}\) and the value at the origin depends on $c$, as shown in \figref{fig:discont_quotient}.
    % \Albert{[Use ``\(\ii\)'' instead of ``\(i\)'' for consistency (I have slight more preference for \(\ii\), since \(i\) is often used as indices (though we don't really have much indices in this article so I guess either is fine).)]}
\end{example}
\begin{example}\label{eg:NonexampleRatio2}
    Let \(\Omega = \RR^2\), \(f(x,y) = x-\ii y\), and \(g(x,y) = x+\ii y\), which have a common simple zero at the origin.  Their quotient is given by \(\varphi(x,y) = \frac{f(x,y)}{g(x,y)} = \frac{x-\ii y}{x+\ii y}\). Just as \autoref{eg:NonexampleRatio1},  this function is also constant along any line passing through the origin but is discontinuous at the origin.
\end{example}

In Examples~\ref{eg:NonexampleRatio1} and \ref{eg:NonexampleRatio2}, the quotient \(f/g\) has a well-defined l'H\^opital limit at the origin when restricted to a one-dimensional path approaching the origin (\eg, along the real axis), as studied in \cite{Carter:1958:HRC}.  However, \(f/g\) is not continuous as a function of two variables.

\vspace{-40pt}
\begin{figure}[h]
    \centering
    % \hspace{20pt}
    % \setlength\unitlength{1pt}
    \begin{minipage}[t]{0.4\textwidth}
    \begin{picture}(210,210)
        \put(0,0)
        {\includegraphics[width=1.0\columnwidth,trim={4cm 0cm 4cm 0cm},clip]{images/eg1_v2.png}}
        % trim={left bottom right top}
        
        % \put(140.8,113.7){\tiny - \( 1.5\)}
        % \put(140.8,63){\tiny - \( 1.0\)}
        % \put(140.8,12.2){\tiny - \( 0.5\)}
      \put(95,12){\small\(x\)}
       \put(35,20){\small\(y\)}
    \end{picture}
    \caption{The real part of the quotient of \(f(x,y)=x+y+\ii y\) and $g(x,y)=x+\ii y$. Blue lines are  level lines, showing that $f/g$ takes a constant value along any line passing through the origin.}
    \label{fig:discont_quotient}
     \end{minipage}
         \hspace{20pt}
          \begin{minipage}[t]{0.4\textwidth}
    \begin{picture}(210,210)
        % \put(0,0){\includegraphics[width=1.0\columnwidth, ,trim={14cm 4cm 14cm 4cm}, clip]{images/eg1_re2D.png}}
        \put(0,0)
        {\includegraphics[width=1.0\columnwidth,trim={4cm 0cm 4cm 0cm},clip]{images/eg2_v2.png}}
        % \put(140.8,113.7){\tiny - \( 2.0\)}
        % \put(140.8,63){\tiny - \( 1.0\)}
        % \put(140.8,12.2){\tiny - \( 0.0\)}
     \put(95,12){\small\(x\)}
       \put(35,20){\small\(y\)}
    \end{picture}
    \caption{The real part of the quotient of 
    % \(f(x,y)=x+x^2+iy\) and $g(x,y)=x+iy$,
     \(f(x,y)= x+x^2+\ii  x+\ii y\) and \(g(x,y) = x + \ii x+\ii y\). The function $f/g$ is continuous but not differentiable at the origin.}
    \label{fig:c0_quotient}
     \end{minipage}
    % \label{fig:projection_onto_sphere}
\end{figure}

\begin{example}\label{eg:c0_quotient}
        % Let \(D \subset \RR^2\) be a small disc around the origin, 
        % \(f(x,y)= x+x^2+\ii y\), and \(g(x,y) = x + \ii y\).
        Let \(\Omega=\RR^2\), \(f(x,y)= x+x^2+\ii  x+\ii y\), and \(g(x,y) = x + \ii x+\ii y\).
        This pair defines a continuous quotient,
        % \(\varphi(x,y) = \frac{f(x,y)}{g(x,y)} = \frac{x+x^2\ii y}{x+\ii y}=1+\frac{x^3-ix^2y}{x^2+y^2}\)
         although it is not differentiable at the origin, as seen in \figref{fig:c0_quotient}. 
\end{example}
This example defines a continuous quotient $f/g$ unlike the previous two examples. We now investigate what makes the difference.

\section{Ratio of derivatives}
Similar to the classical l'H\^opital's rule, the behavior of the quotient of two functions near their common zero is closely related to the quotient of their derivatives. In this section, we explore a more intricate structure that arises near a common simple zero of two complex-valued functions.

Let \(f,g\in C^\infty(\Omega;\CC)\) be smooth complex-valued functions with simple zeros and a common zero set \(\Gamma\).
Then, for all \(\ba\in\Gamma\), the derivatives \(Df|_{\ba}\) and \(Dg|_{\ba}\) are surjective real linear maps to \(\CC\cong\RR^2\).
This defines a unique real matrix \(\bA|_{\ba}\in\RR^{2\times 2}\) (an endormorphism on \(\CC\cong\RR^2\)) that relates the two real linear operators via
\begin{align}
\label{eq:DefinitionOfA}
    Df|_{\ba} = \bA|_{\ba} Dg|_{\ba}.
\end{align}
The matrix representation \(\bA|_{\ba}\) of this relationship depends only on the choice basis for \(\CC\cong\RR^2\), which is taken to be \(1\) and \(\ii\).  In particular, the matrix \(\bA|_{\ba}\) is independent of the choice of coordinate system for \(\Omega\).  

If a change of coordinate \(\Phi\colon\tilde\Omega\to\Omega\) is applied, where \(\Phi(\tilde\ba) = \ba\), and we define \(\tilde f\coloneqq f\circ\Phi\) and \(\tilde g\coloneqq g\circ\Phi\), then the derivatives transform as \(D\tilde f|_{\tilde\ba} = Df|_{\ba} D\Phi|_{\tilde\ba}\) and \(D\tilde g|_{\tilde\ba} = Dg|_{\ba}D\Phi|_{\tilde\ba}\). 
In particular, the relation \(D\tilde f|_{\tilde\ba} = \bA|_{\tilde\ba} D\tilde g|_{\tilde\ba}\) still holds in the transformed coordinates.
That is, the matrix \(\bA\) remains invariant in this change of coordinate. 
\begin{definition}[Ratio of derivatives]
    The function \(\bA\colon\Gamma\to\RR^{2\times 2}\) defined in \eqref{eq:DefinitionOfA} is denoted by
    \begin{align}
    \label{eq:RatioOfDerivatives}
        \left(\frac{Df}{Dg}\right)_{\ba}\coloneqq\bA|_{\ba},\quad \text{for } \ba\in \Gamma.
    \end{align}
\end{definition}

While in general \(f/g\) is discontinuous at \(\Gamma\), as demonstrated in Examples~\ref{eg:NonexampleRatio1} and \ref{eg:NonexampleRatio2}, the nature of the discontinuity is such that the directional limit of \(f/g\) at a point in \(\Gamma\) depends on the path of approach.  This path dependency can be described in terms of the quotient of derivatives \eqref{eq:RatioOfDerivatives}.
\begin{theorem}[Path dependent l'H\^opital rule]\label{thm:PathDependentlHopital}
    Let \(f,g\in C^\infty(\Omega;\CC)\) be smooth complex-valued functions with simple zeros and share a common zero set \(\Gamma\).  
    Suppose \(\gamma\colon (-\varepsilon,\varepsilon)\to\Omega\) is a smooth path such that \(\gamma(0) \in \Gamma\) and \(\gamma'(0)\) is transversal to \(\Gamma\), meaning that \(g_\gamma'\coloneqq \lim_{t\to 0} g(\gamma(t))/t=Dg|_{\gamma(0)}\gamma'(0)\in\CC\) is nonzero.  Then the directional limit of \(f/g\) along \(\gamma\) exists and is given by 
    \begin{align}
    \label{eq:PathDependentlHopital}
        \lim_{t\to 0} \frac{f(\gamma(t))}{g(\gamma(t))} = {1\over g_\gamma'}\left(\frac{Df}{Dg}\right)_{\gamma(0)}g_\gamma'.
    \end{align}
    Here, \(\left(\frac{Df}{Dg}\right)_{\gamma(0)}g_\gamma'\) represents the action of the \(\RR^{2\times 2}\) matrix \({Df\over Dg}\) on \(g_\gamma'\) seen as a real two dimensional vector, while division by  \(g_\gamma'\) is interpreted in terms of complex numbers.  Moreover, the quotient \(f/g\) remains smooth when restricted to the path, meaning that there exists a smooth nowhere-vanishing function \(\varphi\in C^\infty((-\epsilon,\epsilon);\CC)\) such that \(f(\gamma(t)) = g(\gamma(t))\varphi(t)\).
\end{theorem}
We give a proof in \autoref{app:proof_PathDependentlHopital}.
% \Albert{[Need pointer to the proof in appendix]}
Example \autoref{eg:NonexampleRatio1} and \autoref{eg:NonexampleRatio2} in the previous section fall onto the case of \autoref{thm:PathDependentlHopital} and their path-wise limits at the zero is given as in \eqref{eq:PathDependentlHopital}.

\section{Continuous Quotient of Complex Multivariable Functions}
\label{sec:QuotientOfComplexMultivariableFunctions}
The limit \eqref{eq:PathDependentlHopital} for the quotient \(f/g\) of two complex-valued functions with a common simple zero set \(\Gamma\) is generally path-dependent unless the linear operator \(({Df\over Dg})\) satisfies additional conditions.
\begin{definition}[Complex linearity]
\label{def:ConformalEquivalence}
    Two smooth complex-valued functions \(f,g\in C^\infty(\Omega,\CC)\) with simple zeros and a common zero set \(\Gamma\) are said to be \emph{complex linearly related} if the quotient of their derivatives, \(({Df\over Dg})\colon\Gamma\to\RR^{2\times 2}\), is a scaled rotation at every point in \(\Gamma\).  That is, \(({Df\over Dg})\) takes the form
    \begin{align}\label{eq:ScaledRotation}
        \left({Df\over Dg}\right)_{\ba} = 
        \begin{bmatrix}
            u(\ba)&-v(\ba)\\
            v(\ba)&u(\ba)
        \end{bmatrix},\quad\ba\in\Gamma,
    \end{align}
    for some functions \(u,v\colon\Gamma\to\RR\).
\end{definition}

Matrices of the form \eqref{eq:ScaledRotation} can be rewritten as \(
\begin{bsmallmatrix}
    u&-v\\
    v & u
\end{bsmallmatrix}
= r
\begin{bsmallmatrix}
    \cos\theta & -\sin\theta\\
    \sin\theta & \cos\theta
\end{bsmallmatrix}
\) where \(r = \sqrt{u^2 + v^2}\) and \(\theta = \arctan{v\over u}\), showing that the transformation represents a scaled rotation.
Applying the matrix \(\begin{bsmallmatrix}
    u&-v\\
    v & u
\end{bsmallmatrix}\) to an \(\RR^2\) vector is equivalent to complex multiplication by \(u+\ii v = re^{\ii\theta}\).
Another equivalent characterization of a matrix being a scaled rotation is that its two singular values are equal and its determinant is positive.


\begin{theorem}[Continuous quotient of complex functions in arbitrary dimension]\label{thm:ComplexQuotientnD}
    Let \(f,g\in C^1(\Omega;\CC)\) with simple zeros and a common zero set \(\Gamma\).  Then the following statements are equivalent:
    \begin{enumerate}[leftmargin=*,label=(\roman*)]
        \item 
        There exists a non-vanishing continuous function \(\varphi\in C^0(\Omega)\) such that \(f = g\varphi\).       
        \item \(f\) and \(g\) are complex linearly related.
    \end{enumerate}
\end{theorem}
We give a proof in \autoref{app:proof_ ComplexQuotientnD}. The following corollary is an immediate consequence of \autoref{thm:ComplexQuotientnD}.

\begin{corollary}\label{cor:ComplexLinearEquivalenceRelation}
    The notion of complex linearity relation, as defined in \Cref{def:ConformalEquivalence}, constitutes an equivalence relation for functions sharing common simple zeros.
\end{corollary}
\begin{remark}
    Higher regularity of $f$ and $g$ does not, in general, improve the regularity of the quotient in \autoref{thm:ComplexQuotientnD}. Let us consider the pair $f(x,y)=x+x^2+\ii x+\ii y$ and $g(x,y)=x+\ii x + \ii y$ given in Example \ref{eg:c0_quotient}. Its quotient of derivative $Df/Dg$ is an identity matrix at $(0,0)$, hence \autoref{thm:ComplexQuotientnD} applies. But the quotient is merely $C^0$ although both $f$ and $g$ are smooth, as illustrated in \figref{fig:c0_quotient}.
\end{remark}


\subsection{Open problem}
At this point we are unaware of the exact condition for the existence of a smooth quotient. An obvious sufficient condition is that the complex linearity of $Df/Dg$ (\ie, it defines a complex number) extends to an open neighborhood of $\Gamma$. In 2D, this amounts to defining $f$ by multiplying a nowhere vanishing holomorphic function with $g$.

This, of course, is not a necessary condition. For example, consider the pair $f(x,y)=(x+\ii y)e^{\ii(x^2+y^2)},g(x,y)=x+\ii y$, and the pair $f(x,y)=(x+\ii y)(1+x-\ii y)$ (near the origin) and $g(x,y)=x+\ii y$. They both have the  derivative ratios $Df/Dg$ that are complex linear  but do not define a complex number anywhere around the origin. Nevertheless, they both yield smooth quotients. 
We invite the reader to take up the challenge of finding a sharper condition for the existence of a smooth quotient.
% We leave the challenge of finding a sharper condition for the existence of a smooth quotient to future work. 

% \begin{definition}[Conformal equivalence]
% \label{def:conformal_equivalence}    Two smooth complex-valued functions $f,g$ with simple zeros $\Gamma$ and a common zero set is said to be \emph{complex equivalent} if there is an open neighborhood of $\Gamma$ on which the quotient of their derivative $\frac{Df}{Dg}$ is a scaled rotation. 
% \end{definition}

% \begin{definition}[Relative holomorphicity]
% \label{def:relative_holomorphic} Let $g$ be a function . We say $\varphi\in C^\infty(\Omega,\CC)$ is relatively holomorphic to $g$ at $x\in\Omega$ if the matrix-valued quotient $D\varphi/Dg$ defined by
% \begin{align*}
%     D\varphi = \frac{D\varphi}{Dg} Dg
% \end{align*}
% is a complex function in an open neighborhood of $x$.

% independent of slice.

% \end{definition}

% \begin{corollary}\label{cor:EquivalenceRelation}
%     The notion of conformal equivalence, as defined in \Cref{def:ConformalEquivalence}, constitutes an equivalence relation for functions sharing common simple zeros.
% \end{corollary}

% \begin{theorem}[Smooth quotient of complex multivariable functions]\label{thm:SmoothComplexQuotientnD}
%     Let \(f,g\in C^\infty(\Omega;\CC)\) be smooth complex-valued functions with simple zeros and a common zero set \(\Gamma\).  Then the following statements are equivalent:
%     \begin{enumerate}[leftmargin=*,label=(\roman*)]
%         \item 
%         There exists a non-vanishing  function \(\varphi\in C^\infty(\Omega)\) which is relatively hormorphic to $g$ satisfies $f=\varphi g$.
%         \item \(f\) and \(g\) are conformally equivalent.
%     \end{enumerate}
% \end{theorem}

% In fact, Definition \ref{def:conformal_equivalence} is not a sharp condition for the existence of a smooth quotient. For example, $f(x)=z e^{i|z|^2}$ $g(z)=z$ and on a complex plane has a smooth quotient $\varphi(z)=e^{i|z|^2}$, yet the quotient $Df/Dg$ is a complex function only on the zeros $\Gamma$ and it is not so outside the zeros. We may also consider the pair of $g(z)=z$ and $f$ which has only one simple zero $=0$ and is written $f(z)=x+x^2+y^2+iy$ for $x$ near $z=0$. Then $(f,g)$ has a smooth nowhere vanishing quotient $\varphi$,  which is given as the sum of $1$ and a anti holomorphic function: $\varphi(z)=1+x-iy$ near the zero. 


\section{Generalizations}
\label{sec:Generalizations}
The results of Theorems~\ref{thm:RealQuotientnD}, \ref{thm:PathDependentlHopital} and \ref{thm:ComplexQuotientnD} remain valid under any smooth change of coordinates in \(\Omega\in\RR^n\) and do not depend on the Euclidean structure of \(\RR^n\).  Consequently, the domain \(\Omega\) in these theorems can be replaced by any smooth manifold.

Additionally, the assumption of \(C^\infty\)-smoothness for the given functions \(f,g\) in all statements throughout this article can be relaxed to \(C^k\) regularity for any \(k\geq 1\), with the corresponding regularity of \(\varphi\) adjusted to \(C^{k-1}\) accordingly. In fact, the proofs we give in \autoref{app:FirstOrderTaylorExpansion} to \ref{app:proof_ ComplexQuotientnD} are all valid for the general  \(C^k\) regularity case.
 % \Albert{Mention that the proofs in the appendices actually show the general \(C^k\) regularity case.}
%Bibliography

\ifdefined\AMM
\begin{acknowledgment}{Acknowledgments.}
This project was funded in part by the European Research Council (ERC Consolidator
Grant 101045083 CoDiNA) and the National Science Foundation CAREER Award 2239062.
\end{acknowledgment}
\else
% \subsection*{Acknowledgments}
\begin{acknowledgment}{Acknowledgments.}
This project was funded in part by the European Research Council (ERC Consolidator
Grant 101045083 CoDiNA) and the National Science Foundation CAREER Award 2239062.
\end{acknowledgment}
\fi

\bibliographystyle{vancouver}
\bibliography{Reference}
\appendix

\section{First Order Taylor Expansion}
\label{app:FirstOrderTaylorExpansion}
The smoothness statements in Theorems~\ref{thm:RealQuotient1D}--\ref{thm:PathDependentlHopital} follow from the regularity of the remainder terms in the Taylor expansions.

Suppose \(f\colon I\to\RR\) is a continuously differentiable function of one real variable over an open interval \(I\ni a\).  Then \(f\) can be expressed using the first-order Taylor formula as \(f(x) = f(a) + f'(a)(x-a) + (x-a)h(x)\), where the function \(h\) satisfies \(\lim_{x\to a}h(x) = 0\).  This form of the Taylor remainder is known as \emph{Peano's form}, first characterized by 
Giuseppe Peano in 1889 \cite{Peano:1889:UNF,Persson:2017:HST}.   

If we further assume that \(f\in C^\infty(I)\), then the Peano remainder function \(h\) is also \(C^\infty\)-smooth.  This regularity result is known as Hadamard's Lemma, which states that \(f(x) = f(a) + (x-a)g(x)\) for some \(g\in C^\infty(I)\).  An elementary proof of Hadamard's lemma can be found in \cite{Nestruev:2003:SMO}.

More generally, if \(f\in C^k(I)\) for some \(k\geq 1\), then the first-order Peano remainder function \(h(x)\) (or \(g(x) = h(x)+f'(a)\) in the form of Hadamard's lemma) belongs to \(C^{k-1}(I)\).  This was shown by Hassler Whitney in 1943 \cite{Whitney:1943:DRT}.

\begin{lemma}[Whitney 1943 \cite{Whitney:1943:DRT}]
\label{lemma:Whitney}
    Suppose \(f\in C^k(I)\) with \(k\geq 1\) on an open interval \(I\ni 0\).  Then there exists a function \(g\in C^{k-1}(I)\) such that
    \(
        f(x) = f(0) + xg(x),
    \)     
    with \(g^{(j)}(0) = {1\over j+1}f^{(j+1)}(0)\) for \(0\leq j\leq k-1\).
\end{lemma}
\begin{proof}
    For \(x\neq 0\), the finite difference quotient \(g = {f(x) - f(0)\over x}\) can be written as \(g(x) = {1\over x}\int_0^xf'(t)\, dt\).   The derivatives of \(g(x)\) admit the following integral formula:
    \begin{align*}
    % \label{eq:WhitneyMagicFormula}
        {d^j\over dx^j}g(x) = {1\over x^{j+1}}\int_0^x t^j f^{(j+1)}(t)\, dt,\quad j=0,\ldots,k-1.
    \end{align*}
    This identity can be verified by repeated integration by parts and mathematical induction; we omit the details here.  
    Since \(f\in C^k(I)\), it follows that \(g\in C^k(I\setminus\{0\})\subset C^{k-1}(I\setminus\{0\})\).
    It remains to show that \(g^{(j)}(x)\) has a limit as \(x\to 0\), for \(0\leq j\leq k-1\).
    Define \(L_j \coloneqq {1\over j+1}f^{(j+1)}(0)\), which can also be expressed as \(L_j = {1\over x^{j+1}}\int_0^x t^jf^{(j+1)}(0)\, dt \) for any \(x\neq 0\).  Then we compute:
    \begin{align*}
        \left|g^{(j)}(x) - L\right| & = 
        \left|{1\over x^{j+1}}\int_0^x t^j\left(f^{(j+1)}(t) - f^{(j+1)}(0)\right)\, dt\right|\\
        &\leq \underbrace{{1\over x^{j+1}}\int_0^x t^j\, dt}_{{1\over j+1}} \max_{t\in [0,x]}\left|f^{(j+1)}(t) - f^{(j+1)}(0)\right|.
    \end{align*}
    Since \(f^{(j+1)}\) is continuous at \(0\), this expression tends to zero as \(x\to 0\).  Therefore, each \(g^{(j)}\) extends continuously to \(x=0\), with \(g^{(j)}(0) = L_j\).  Hence, \(g\in C^{k-1}(I)\).
\end{proof}

Whitney also extended the result of \Cref{lemma:Whitney} to multidimensions in \cite{Whitney:1943:DRT}.  Suppose \(f\colon\RR^n\to\RR\) is defined in a neighborhood of the origin.  Then by fixing each \(x_2,\ldots,x_n\) and applying the same argument as \Cref{lemma:Whitney} on the variable \(x_1\), we obtain
\begin{align}
\label{eq:WhitneyTowardsMultidiemsion}
        f(x_1,x_2,\ldots,x_n) = f(0,x_2,\ldots,x_n) + x_1 g_1(x_1,x_2,\ldots,x_n)
\end{align}
where \(f(0,x_2,\ldots,x_n)\) is still \(C^k\) and \(g_1\in C^{k-1}\).  The continuity of \({\partial^{i_1+\cdots+i_n}\over\partial x_1^{i_1}\cdots \partial x_n^{i_n}}g_1\) for \(i_1+\cdots+i_n\leq{k-1}\) is the result of applying \Cref{lemma:Whitney} on the function \({\partial^{i_2+\cdots +i_n}\over\partial x_2^{i_2}\cdots\partial x_n^{i_n}}f\).

Repeating the process of \eqref{eq:WhitneyTowardsMultidiemsion} yields
\begin{align}
\label{eq:WhitneyMultidimensionScalar}
        f(\bx) = f(\bzero) + x_1g_1(x_1,\ldots,x_n) + x_2g_2(x_2,\ldots,x_n)+\cdots+ x_n g_n(x_n)
\end{align}
where \(g_i\in C^{k-1}\) and \(g_i(\bzero) = {\partial f\over\partial x_i}(\bzero)\).

The result \eqref{eq:WhitneyMultidimensionScalar} also holds for vector-valued functions \(\bff = \begin{bsmallmatrix}
    f_1\\|\\f_m
\end{bsmallmatrix}\colon\RR^n\to\RR^m\) by applying \eqref{eq:WhitneyMultidimensionScalar} component-wise, which is concluded by the following lemma.
In the lemma we single out the factors \(\sum_{i}{\partial\bff\over\partial x_i}(\bzero)x_i = D\bff(\bzero)\bx\) to resemble Peano's form of Taylor expansion.

\begin{lemma}[First-order Taylor formula with regularity]
\label{lemma:Taylor}
Let \(\bff\colon U\subset\RR^n\to\RR^m\) be a \(C^k\)-continuous vector-valued function on a neighborhood \(U\) of a given point $\ba\in \RR^n$.
Then there exists \(\bH\in C^{k-1}(U;\RR^{m\times n})\) matrix-valued function so that \(\bH(\ba) = \bzero\) and
\begin{align*}
    \bff(\bx) = \bff(\ba) + D\bff(\ba)(\bx-\ba) + \bH(\bx)(\bx-\ba),\quad\bx\in U.
\end{align*}
\end{lemma}

\section{Proof of Theorem \ref{thm:RealQuotient1D}}
\label{app:ProofOfRealQuotient1D}



% \begin{lemma}
%     Let \(f\colon I\to\FF\), with \(\FF = \RR\) or \(\CC\), be a \(C^k\)-smooth function of one variable over a real open interval \(I\), with \(k\geq 1\).  Suppose \(f(a) = 0\) at \(a\in I\) and \(f'(a) \neq 0\). Then there exists a neighborhood \(U\subset I\) of \(a\) where \(f(x) = f'(a)(x-a)F(x)\) for some \(C^{k-1}\)-smooth function \(F\colon U\to\FF\) satisfying \(F(a) = 1\) and \(F(x)\) is non-vanishing on \(U\).
% \end{lemma}
% \begin{proof}
%     By the Taylor Theorem with Peano's remainder, there exists a function \(h(x)\in C^{k-1}(I;\FF)\) such that \(f(x) = f'(a)(x-a) + (x-a)h(x)\) with \(\lim_{x\to a}h(x) = 0\).  Note that the classical Taylor's Theorem which applies to real-valued function extends to complex-valued function by applying Taylor's Theorem on real and imaginary part separately.  Taking the common factor \(f'(a)(x-a)\) out to obtain \(f(x) = f'(a)(x-a)(1+h(x)/f'(a))\) and define \(F(x) = (1+h(x)/f'(a))\in C^{k-1}(I;\FF)\).  Therefore \(f(x) = f'(a)(x-a)F(x)\).  The limit of \(F(x)\) at \(x\to a\) is 1.  By continuity of \(F\) we have \(F(a)=1\) and there exists a nieghborhood of \(a\) in which \(F\) is nonzero.
% \end{proof}
\begin{proof}[Proof of \autoref{thm:RealQuotient1D}]
    Let \(\Gamma\) be the zero set of \(f\in C^k(I)\) and \(g\in C^k(I)\).  For \(x\in I\setminus\Gamma\) we define \(\varphi\) by \(\varphi(x) = f(x)/g(x)\), which belongs to \(C^{k}(I\setminus\Gamma)\subset C^{k-1}(I\setminus\Gamma)\) and non-vanishing.  It remains to show that \(\varphi\) is \(C^{k-1}\)-smoothly defined in a neighborhood of \(\Gamma\).  
    By \Cref{lemma:Whitney}, near a simple zero \(a \in \Gamma\) of both \(f\) and \(g\), we can express  
\[
f(x) = f'(a)(x-a)F(x), \quad g(x) = g'(a)(x-a)G(x)
\]
for \(x\) in a neighborhood \(U\) of \(a\), where \(F,G\in C^{k-1}(U)\), \(F(a) = G(a) = 1\), and both \(F(x)\) and \(G(x)\) are nonzero on \(U\). 
Thus, for \(x \in U\), we define  
\[
\varphi(x) \coloneqq \frac{f'(a)}{g'(a)} \cdot \frac{F(x)}{G(x)},
\]
which is smooth, non-vanishing,  $\varphi(a)=f'(a)/g'(a)$, and agrees with \(f(x)/g(x)\) on \(U\setminus\{a\}\).  
\end{proof}

\section{Proof of Theorem \ref{thm:RealQuotientnD}}
\label{app:ProofOfRealQuotientnD}
\begin{proof}[Proof of \autoref{thm:RealQuotientnD}]
    Let \(\Sigma\) be the common zero set of \(f,g\in C^k(\Omega)\).  On \(\Omega\setminus\Sigma\) define \(\bx\in\Omega\setminus\Sigma\) define \(\varphi(\bx) = f(\bx)/g(\bx)\), which belongs to \(C^{k}(\Omega\setminus\Sigma)\) and is non-vanishing.
    It remains to show that \(\varphi\) is \(C^{k-1}\)-continuously defined in a neighborhood of any point of \(\Sigma\).  Let \(\ba\in\Sigma\).  By the implicit function theorem, there exists a neighborhood \(U\) of \(\ba\), a neighborhood \(\tilde U\subset\RR^n\) of the origin \(\bzero\in\RR^n\), and a \(C^k\) bijection \(\Phi\colon \tilde U\xrightarrow{\simeq} U\) such that \(\Phi(\bzero) = \ba\) and \(\tilde\Sigma = \Phi^{-1}(\Sigma)\) is the coordinate hyperplane \(\tilde\Sigma = \{(x_1,\ldots, x_n)\,|\,x_1 = 0\}\cap\tilde U\).
    Define \(\tilde f = f\circ\Phi\) and \(\tilde g = g\circ\Phi\), both of which are \(C^k\) functions on \(\tilde U\) and have simple zeros on \(\tilde\Sigma\).  
    In this coordinate system, \(\nabla \tilde f(0,x_2,\ldots,x_n) = {\partial \tilde f\over\partial x_1}(0,x_2,\ldots,x_n)\be_1\) and \(\nabla \tilde g(0,x_2,\ldots,x_n) = {\partial \tilde g\over\partial x_1}(0,x_2,\ldots,x_n)\be_1\), where \(\be_1\) is the coordinate vector \((1,0,\ldots,0)\), and \({\partial\tilde f\over\partial x_1}\) and \({\partial\tilde g\over\partial x_1}\) are nonvanishing \(C^{k-1}\) function on \(\tilde\Sigma\).
    By \eqref{eq:WhitneyTowardsMultidiemsion}, there exists \(p,q\in C^{k-1}(\tilde U)\) such that for \((x_1,\ldots,x_n)\in \tilde U\)
    \begin{align*}
        \tilde f(x_1,\ldots,x_n) = x_1 p(x_1,\ldots,x_n),\quad\tilde g(x_1,\ldots,x_n) = x_1q(x_1,\ldots,x_n)
    \end{align*}
    and \(p = {\partial\tilde f\over\partial x_1}\neq 0\) and \(q = {\partial\tilde g\over\partial x_1}\neq 0\) on \(\tilde\Sigma\).
    Therefore there exists a neighborhood of \(\tilde\Sigma\), without loss of generality \(\tilde U\), on which \(p\) and \(q\) are nonvanishing.  Therefore, \({\tilde f\over\tilde g} = {p\over q}\) is a nonvanishing \(C^{k-1}\) continuous function on \(\tilde U\) and agrees with \(\varphi\circ \Phi\) on \(\tilde U\).  Thus \(\varphi\) extends to a \(C^{k-1}\)-continuous function in a neighborhood of \(\ba\).
\end{proof}

\section{Proof of Theorem \ref{thm:PathDependentlHopital}}\label{app:proof_PathDependentlHopital}

To prove \Cref{thm:PathDependentlHopital} we generalize \Cref{lemma:Whitney} and \Cref{thm:RealQuotient1D} that apply to complex-valued functions.

% To prove the theorem, we use a complex version of Theorem \ref{thm:RealQuotient1D}, which can be obtained with the complex  version of Lemma 1. For completeness, we state them as lemmas. 
% \Albert{[``complex version'' sounds pretty vague.  We can shorten this intro to ``To prove \Cref{thm:PathDependentlHopital} we generalize \Cref{lemma:Whitney} and \Cref{thm:RealQuotient1D} that apply to complex-valued functions.'']}

\begin{lemma}
\label{lemma:ComplexWhitney}
    Suppose \(f\in C^k(I,\CC)\) with \(k\geq 1\) on an open interval \(I\ni 0\).  Then there exists a function \(g\in C^{k-1}(I,\CC)\)
    % \Albert{[shouldn't this also be complex-valued?]} such that
    \(
        f(x) = f(0) + xg(x),
    \)     
    with \(g^{(j)}(0) = {1\over j+1}f^{(j+1)}(0)\) for \(0\leq j\leq k-1\).
\end{lemma}
\begin{proof}
We apply \Cref{lemma:Whitney} to the real and imaginary parts of \(f\) to obtain the real and imaginary parts of \(g\) holding the stated properties.
% \Albert{We apply \Cref{lemma:Whitney} to the real and imaginary parts of \(f\) to obtain the real and imaginary parts of \(g\).}
\end{proof}

\begin{lemma}[Smooth quotient of complex functions on a real line]\label{lem:ComplexQuotient1D}
     Let \(f,g\in C^\infty(I,\CC)\) be smooth complex-valued functions on a real interval \(I\). Suppose \(f,g\) share common zeros, and their derivatives $f'$ and $g'$ do not vanish on the zero set. Then there exists a non-vanishing smooth complex-valued function \(\varphi\in C^\infty(I,\CC)\) such that \(f = g\varphi\).
\end{lemma}
\begin{proof}
    The proof follows the same steps as the proof of \Cref{thm:RealQuotient1D} (\Cref{app:ProofOfRealQuotient1D}), with \Cref{lemma:Whitney} replaced by \Cref{lemma:ComplexWhitney}.
\end{proof}

\begin{proof}[Proof of Theorem \ref{thm:PathDependentlHopital}]
    Consider smooth functions $f_\gamma(t)\coloneqq f(\gamma(t))$ and $f_\gamma(t)=g(\gamma(t))$ on the interval $(-\varepsilon,\varepsilon)$. 
    By \Cref{lem:ComplexQuotient1D},
    the function $\varphi\coloneqq  f_\gamma/ g_\gamma$ is smooth and nowhere vanishing, with $\varphi(0)=f'_\gamma(0)/g'_\gamma(0)$.
    Noting that $g_\gamma'(0)=Dg|_{\gamma(0)}\gamma'(0)$ and $f_\gamma'(0)=Df|_{\gamma(0)}\gamma'(0)= \frac{Df}{Dg} g_\gamma'(0)$, we obtain the stated expression.
\end{proof}

\section{Proof of Theorem \ref{thm:ComplexQuotientnD}}\label{app:proof_ ComplexQuotientnD}

% \begin{proof}[Proof of \autoref{thm:ComplexQuotientnD}]
%     We first show that $(ii)$ implies $(i)$. The continuity of $\varphi=f/g$ on $\Omega\setminus f^{-1}(0)$ is trivial. We prove now $\varphi$ is defined and continuous on at each point $p$ of $f^{-1}(0)$. In a small open neighborhood $B$ of $p$ we take a local coordinate $(\bx, \by)=(x_1,x_2, y_1,\ldots, y_{n-2})$ where $f^{-1}(0)\subset B=\{(\bx,\by)\in B \,\mid\, \bx=\mathbf{0}\}$. 

%     Note that for each $\by$, consider the slice $B_\by$ of $B$. On $B_\by$ 

% \TODO{To be continued.}

% And $\varphi(\bzero,\by)$ is continuous in $\by$ since $\varphi(\bzero,\by)=Df/Dg(\bzero,\by)\in \CC$ for each $\by$ and the matrix-valued function $Df/Dg$ is smooth everywhere by assumption.

% Now $\varphi$ may appear to be continuous at $(\bzero,\by)$ as $\varphi(\bzero,\by)$ is continuous in $\by$, and for each curve $\gamma$ transversally intersecting with $\{\bx=0\}$ at $\gamma(0)=(\bzero,\by)$, $\varphi(\gamma(s))$ is continuous and $\varphi(\gamma(0))=\varphi(\bzero,\by)$. This however does not exclude the possibility of the existence of a path $\gamma$ that tangentially touches $\{\bx=0\}$ at $(\bzero,\by)=\gamma(0)$ and $\lim_{s\to 0}\varphi(\gamma(s))\neq \varphi(\bzero,\by)$.


% \TODO{To be continued.}

%     We want to use the triangle inequality
%     \begin{align}
%         |\varphi(\bx, \by)-\varphi(\bzero,\bzero)|\leq |\varphi(\bzero,\by)-\varphi(\bzero,\bzero)|+|\varphi(\bzero,\by)-\varphi(\bx,\by)|.
%     \end{align}
%     The control of the first term in the right-handed-side is easy since $\varphi(\bzero,\by)$ is continuous in $\by$. However, controlling the second term is non-trivial. For each given $\epsilon>0$, there exists $\delta_\by>0$ for which $|\bx|<\delta_\by$ implies $|\varphi(\bzero,\by)-\varphi(\bx,\by)|<\epsilon$. However, the strict positivity $\inf_{\by\in I}\delta_\by >0$ on some small domain $I=[-\by_0,\by_0]^{n-2}$ is unclear at this point. Namely, we have to show the equicontinuity of $M\coloneqq\{\varphi_\by\}_{\by\in I}$. 

%     To make our setting applicable for the (arguably) converse direction of the Ascoli-Arzela theorem,
%     we consider the restriction of each $\varphi_\by$ on a small closed box $X=[-x_0,x_0]^2$ so that they are bounded and belong to Banach space
%     $(C(X),\|\cdot\|_{\sup})$. For the restricted functions, we still use the same notation $\varphi_\by$ for simplicity. 

%     We now show that $M$ is compact in $C(X)$. For any sequence $\{\varphi_{n}\}_n$ in $M$ is of the form $\{\varphi_{\by_n}\}_n$ and there is a subsequence $\by_{n_j}$ converging to some $y\in I$ as $I$ is compact. Then $\lim_{j->\infty}\|\varphi_{\by_{n_j}}\to \varphi_\by\|_{\sup}$ as  $\varphi_{\by_{n_j}}(\bx)\to \varphi_\by(\bx)$ for $\bx\neq \bzero$ due to the continuity of $\varphi$ on $\{\by \neq \bzero\}$ and for $\bx=\bzero$ due to the continuity of $\varphi(\bzero,\by)$ in $\by$.
    
%     there is a convergent subsequence $f_{y_{n_j}}\to f_y$ for some $y$ 
    
% \end{proof}

% \Sadashige{I commented out my long proof below.}

% The Peano reminder $R(\bx)\coloneqq H(\bx)\cdot \bx$ in Lemma 2 is in $o(|\bx|)$. That is, 
% \begin{align}
%     \lim_{\bx\to \bzero} \frac{R(\bx)}{|\bx|}=0
% \end{align}
% where $| \bx|=\sqrt{\sum_i x_i^2}$ is the standard Euclidean norm.

% In our case, we have a stronger result. Let us assume the setting of Theorem 4. Then around each point of the zero set $\Gamma$, we may take without loss of generality a local coordinate system $\bx=(\bx^\perp, \bx^\parallel)$ with $2$ dimensional $\bx^\perp$ and $n-2$ dimensional $\bx^\parallel$ so that  $\Gamma$ is identical with the local $n-2$ dimensional hyperplane $\{\bx=(\bx^\perp,\bx^\parallel) \mid\, \bx^\perp=\bzero\}$. Here and in the rest of the appendix, the dimension of the zero vector $\bzero$ depends on the context. Then, we have the following:
% \begin{lemma}[Refinement of Lemma \ref{lemma:Taylor}]
% In the above local coordinates, we have $R\in o(|\bx^\perp|)$, that is,
% \begin{align}
%     \lim_{\bx\to \bzero} \frac{R(\bx)}{|\bx^\perp|}=0.
% \end{align}
% \end{lemma}

% \begin{proof}
% In our local coordinates, we have
% \[
% Df(\bx)=
% \begin{pmatrix}
% \begin{array}{ccccc}
% \multicolumn{2}{c}{\multirow{2}{*}{\scalebox{1.2}{$DF$}}}
%   & \multicolumn{3}{c}{\multirow{2}{*}{\scalebox{1.2}{$\ RF\ $}}} \\
%  &  &  &  &  \\
% \end{array}
% \end{pmatrix}
% \begin{pmatrix}
% \begin{array}{cc}
% \bx^\perp\\
% \bx^\parallel
% \end{array}
% \end{pmatrix}
% \]
% where $DF$ is a $2$ by $2$ matrix valued function, which is invertible on the simple zeros $\Gamma$, and $RF$ is a $2$ by $(n-2)$ matrix valued function, which vanishes on $\Gamma$.

% Lemma \ref{lemma:Taylor} states
% \begin{align}
%    f(\bx)=f(\bzero)+Df(\bzero)\cdot \bx + H(\bx)\cdot \bx
% \end{align}
% where $H(\bx)\to \bzero$ as $\bx\to \bzero$. Hence we have for $R(\bx)=H(\bx)\cdot \bx$ that, 
% \begin{align}
%     R(\bx) &= f(\bx)- f(\bzero)- Df(\bzero)\cdot (\bx)\\
%     &= f(\bx^\perp,\bx^\parallel) - f(\bzero,\bx^\parallel) - DF(\bzero)\cdot \bx^\perp\\
%     &=\left( DF(t_{\bx} \bx^\perp,\bx^\parallel) - DF(\bzero) \right)\cdot \bx^\perp
% \end{align}
% where $t_{\bx}\in [0,1]$. Here we used $f(\bzero)=f(\bzero,\bx^\parallel)=0$ and applied the mean value theorem to the function $\tau:[0,1]\to \CC$ given by $\tau(t)\coloneqq f(t \bx^\perp,\bx^\parallel)$. 

% Since $DF(t_{\bx}\bx^\perp,\bx^\parallel)\to DF(\bzero)$ as $\bx\to \bzero$, we have
% \begin{align}
%     \lim_{\bx\to \bzero}\frac{R(\bx)}{|\bx^\perp|}
%     \leq  \lim_{\bx\to \bzero}\frac{\| DF(t_\bx \bx^\perp,\bx^\parallel) - DF(\bzero,\bx^\parallel) \|\,|\bx^\perp|}{|\bx^\perp|}\to 0
% \end{align}
% where $\|\cdot \|$ is the operator norm defined by $\|A\|=\sup_{|x|=1}|Ax|$ for a linear operator $A:\RR^2\to \CC\simeq \RR^2$. 
% \end{proof}

On a 2D domain, \autoref{thm:ComplexQuotientnD} is a direct consequence of \autoref{thm:PathDependentlHopital}. 
If \(f\) and \(g\) are complex linearly related, 
with  the quotient of derivative
% in \eqref{eq:PathDependentlHopital} is a scaled rotation 
\(({Df\over Dg}) = \begin{bsmallmatrix}
    u&-v\\
    v & u
\end{bsmallmatrix}\) with some $u,v\in \RR$ at each zero, 
then \eqref{eq:PathDependentlHopital} simplifies to 
\begin{align*}
    \lim_{t=0}{f(\gamma(t))\over g(\gamma(t))} = {1\over g'}(u+\ii v)g' = u+\ii v
\end{align*}
which is a single complex number that depends only on the point \(\gamma(0)\) and is independent of the path direction \(\gamma'\). This gives a proof of \autoref{thm:ComplexQuotientnD} in 2D.
% \Albert{[I'm not sure if I understand this.  Why can the proof only consider \(f,g\) complex linear?  Do we mean they are complex linearly related?]}

In higher dimensions, however, \autoref{thm:ComplexQuotientnD} does not automatically follow from \Cref{thm:PathDependentlHopital}.
Let us consider any 2D plane which transversely intersects $\Gamma$, on which the situation reduces to the 2D setting discussed above. Hence the quotient function $\varphi$ is continuous along any path that intersects $\Gamma$ transversely. In addition, $\varphi$ is smooth when restricted either to $\varphi|_{\Gamma}$ or to $\varphi_{\Omega\setminus\Gamma}$ by construction. These properties, however, are not sufficient to conclude the continuity of $\varphi$. In fact, we can construct a function that satisfies all these properties that is still undefined on $\Gamma$, as the following example illustrates. 
\newline
% 
% \vspace{10pt}
\begingroup
% \begin{wrapfigure}{r}{0.15\textwidth}
% \centering
% % trim={left bottom right top}
% \includegraphics[trim={2cm 0 2cm 0},clip, width=0.12\textwidth]{images/bump_functs.png}
% \caption{test.}
% \label{fig:bump_functions}
% \end{wrapfigure}
% Adjusted settings
\setlength{\intextsep}{10pt}%
\setlength{\columnsep}{10pt}%
% \vspace{3cm}
\begin{wrapfigure}{r}{0.2\textwidth}
% % trim={left bottom right top}
  % \centering\includegraphics[width=1.0\linewidth,trim={13cm 1cm 13cm 0}, clip]{images/bump_functions_no_labels.png}
  % \caption{Basic layout}
      \begin{picture}(0,150)
        \put(0,0){\includegraphics[width=.2\columnwidth,trim={13cm 1cm 13cm 0}, clip]{images/bump_functions_no_labels.png}}
        \put(45,22){\small\(x\)}
        \put(16,80){\small\(z\)}
    \end{picture}
    % \caption{ Cut-off view of $xz$-plane. Discs are the supports of the bump functions $b_n$, lying on the curve $z=2\sqrt{x}$. }
\end{wrapfigure}
\begin{example}\label{eg:infinite_bump_functions}
    Consider a collection of smooth bump functions $\{b_n\}_{n\in \NN}$ on $\RR^3$ such that each $b_n$ takes its maximum value $1$ at the center $(2^{-n},0, 2^{1-{n\over 2}})$ of its support, given as an open ball of radius $2^{-n}$. Their centers lie on the curve $z=2\sqrt{x}$, and the supports touch $z$-axis, as visualized in the right image.
    Using these bump functions, let us define $\varphi(x,y,z)= \sum_{n=0}^\infty b_n(x,y,z)$. 
    Notice that any path intersecting $z$-axis transversely passes through only finitely many of the above balls. Hence, along such paths $\gamma:(-\epsilon,\epsilon)\to \RR^3$ with $\gamma(0) = \mathbf{0}$, we have $\lim_{t\to0} \varphi(\gamma(t))=0$. However, any curve $(t,0,a \sqrt{t})$ with $a>2$, intersects these balls infinitely many times as $t\to 0$, making $\varphi$ discontinuous at the origin.
    % \Albert{I don't think it is clear to me how this example fits the overall argument.  All this example is saying is that seeing a function with consistent limit along straight lines is not enough to conclude continuity. But it's not clear from the paragraph before the example that we are in this situation.  So the argument in the middle of the example that suddenly brings in ``straight line'' is quite confusing.}
\end{example}
 \endgroup
 We now prove \autoref{thm:ComplexQuotientnD} in arbitrary dimension. The key component is the simple zero condition of $f$ and $g$. It ensures that the leading order of these functions on $\Omega\setminus \Gamma$ consists of the coordinates perpendicular to $\Gamma$. This results in the continuity of the quotient along any path approaching to $\Gamma$, possibly non-transversely, as in Example \ref{eg:infinite_bump_functions}.

\begin{proof}[Proof of \autoref{thm:ComplexQuotientnD}]
{$(i)\to (ii)$.} 
It follows from the $C^1$ regularity of $g$ and the continuity of $\varphi$ that $Df(\ba)=\varphi(\ba)Dg(\ba)$ at each zero $\ba$. This shows that \(Df/Dg = \begin{bsmallmatrix}
    \Re\varphi(\ba)&-\Im\varphi(\ba)\\\Im\varphi(\ba)&\Re\varphi(\ba)
\end{bsmallmatrix}\) and therefore \(f\) and \(g\) are complex linearly related.

% We show that $Df/Dg(\ba)$ is a complex number \Albert{[(more like a matrix of a particular form that can be viewed as a complex number) maybe a pointer to that matrix representation]} at any $\ba\in \Gamma$.
% It follows from the $C^1$ regularity of $g$ and the continuity of $\varphi$ that $Df(\ba)=\varphi(\ba)Dg(\ba)$ even if $\varphi$ is not differentiable.
% \Albert{[Not sure we need to say ``even if \(\varphi\) is not differentiable.'']}
% % 
% \suggest{
% To see this, let us recall that $D f(\ba)$ is defined as a unique linear operator $L:\RR^n\to \CC$ such that 
% \begin{align*}\label{eq:def_of_Df}
%    \lim_{\bx \to \ba} \frac{f(\bx)-f(\ba) - L(\bx-\ba) }{|\bx-\ba|}=0.
% \end{align*}
% Since $g(\bx)=Dg(\ba)(\bx-\ba)+H(\bx)(\bx-\ba)$ for some continuous function $H$ such that $H(\ba)=0$ due to \autoref{lemma:Taylor}, we have
% \begin{align*}\label{eq:def_of_Df}
%    \frac{f(\bx)-f(\ba) - L(\bx) }{|\bx|}
%    &= \frac{|\varphi(\bx)(Dg(\ba)+H(\bx-\ba))(\bx-\ba)-L(\bx-\ba)|}{|\bx-\ba|}\\
%    &\leq \|\varphi(x)(Dg(\ba)+H(\bx-\ba))-L\|
% \end{align*}
% where $\|\cdot\|$ denotes the operator norm. This converges to $0$ as $\bx\to \ba$ if $L=\varphi(\ba)Dg(\ba)$.
% }\Sadashige{I think we may drop the above part as it is almost trivial.}
% \Albert{I agree.  Just state as if it is straightforward.  \(Df(\ba) = \varphi(\ba)Dg(\ba)\) implies that \(Df/Dg = \begin{bsmallmatrix}
%     \Re\varphi(\ba)&-\Im\varphi(\ba)\\\Im\varphi(\ba)&\Re\varphi(\ba)
% \end{bsmallmatrix}\) and therefore \(f\) and \(g\) are complex linearly related.}

\vspace{5pt}
\noindent $(ii)\to (i)$.
    Let \(\Gamma\) be the common zero set of \(f,g\in C^1(\Omega;\CC)\).  
    The quotient \(\varphi(\bx) = f(\bx)/g(\bx)\) for \(\bx\in \Omega\setminus\Gamma\), is \(C^1\)-continuous and nonvanishing.
    It remains to show that \(f/g\) extends to a \(C^{0}\)-function in a neighborhood of any point \(\ba\in \Gamma\).  

    % As in the proof of \autoref{thm:RealQuotientnD}, we may take without loss of generality a local coordinate system centered at $\ba$ and $\Gamma$ is 
    % we may take without loss of generality a local coordinate system $\bx=(\bx^\perp, \bx^\parallel)$ centered at $\ba$ with $2$ dimensional $\bx^\perp$ and $n-2$ dimensional $\bx^\parallel$ so that  $\Gamma$ is identical with the local $n-2$ dimensional hyperplane $\{\bx=(\bx^\perp,\bx^\parallel) \mid\, \bx^\perp=\bzero\}$.

    
    As in the proof of \autoref{thm:RealQuotientnD}, the implicit function theorem assures the existence of a neighborhood \(U\) of \(\ba\), a neighborhood \(\tilde U\) of the origin \(\bzero\in\RR^n\), and a \(C^1\) bijection \(\Phi\colon\tilde U\xrightarrow{\simeq} U\) such that \(\Phi(\bzero) = \ba\) and \(\tilde\Gamma = \Phi^{-1}(\Gamma)\) is the subspace \(\tilde\Gamma = \{\bx=(x_1,x_2,\ldots,x_d)\,|\,  x_1=x_2=0\}\cap\tilde U\). Let us write $\bx^\perp=(x_1,x_2)$ and $\bx^\parallel=(x_3,\ldots,x_n)$ for simplicity.
    % \Albert{[readability of the last bit here still needs some work. Just first characterize \(\Phi\) clearly, with maximal clarity when talking about \(\tilde\Gamma\).  Then, separately, define new notations \(\bx^\parallel,\bx^\bot\) for shorthand. ]}
    Now let \(\tilde f = f\circ\Phi\) and \(\tilde g = g\circ\Phi\), which are \(C^1\) functions vanishing on \(\tilde\Gamma\).

\begin{claim}\label{cl:P and Q}
There exists \(\bP,\bQ\in C^{0}(\tilde U;\RR^{2\times 2})\) such that 
    \begin{align}\label{eq:P and Q}
        \tilde f(\bx) = \bP(\bx)\bx^\perp,\quad
        \tilde g(\bx) = 
        \bQ(\bx)\bx^\perp,
    \end{align}
    with $\bP(\bzero,\bx^\parallel)=\nabla_{\bx^\perp}\tilde f(\bzero,\bx^\parallel)$ and $\bQ(\bzero,\bx^\parallel)=\nabla_{\bx^\perp}\tilde g(\bzero,\bx^\parallel)$ being invertible, where the equalities  in \eqref{eq:P and Q} are the natural identification between $\CC\cong \RR^2$. 
\end{claim}
We will prove this claim at the end.
Then from the continuity of $\bP,\bQ$ and the assumption that \(f\) and \(g\) are complex linearly related, it follows that there exists a continuous complex-valued function \( \lambda\) on $\tilde \Gamma$ such that
% \(Df = \lambda Dg\) and
% \(D\tilde f = \tilde\lambda D\tilde g\) where \(\tilde \lambda = \lambda\circ\Phi\). 
     \(\bP = \Lambda\bQ\) on \(\tilde\Gamma\) with the matrix-valued function  $\Lambda=\begin{bsmallmatrix}
         \Re\lambda & -\Im\lambda\\\Im\lambda&\Re\lambda
     \end{bsmallmatrix}$.
     % \Albert{[This \(\tilde\lambda\) is converted to a 2x2 matrix.  Maybe let \([\lambda]\) denote the matrix \(\begin{bsmallmatrix}\Re\lambda & -\Im\lambda\\\Im\lambda&\Re\lambda\end{bsmallmatrix}\) for a complex \(\lambda\in\CC\).  Then you can say there exists \(\tilde\lambda\in C^0(\tilde\Gamma,\CC)\)] so that \(\bP = [\tilde\lambda]\bQ\).}
% Note that $\tilde \lambda= \bP \bQ^{-1}
% =\nabla_{\bx^\perp}\tilde f(\bzero,\bx^\parallel) \nabla_{\bx^\perp}\tilde g(\bzero,\bx^\parallel)^{-1}
% $ is continuous since $\bP, \bQ$ are continuous on $\tilde \Gamma$.
    % 
Using this, we define a function $\tilde \varphi$ on $\tilde \Gamma$ by
    \begin{align*}
    \tilde\varphi(\bx)=
    \begin{cases}
			\tilde f(\bx)/\tilde g(\bx), & \bx \notin \tilde\Gamma,\\
            \lambda(\bx^\parallel), &  \bx \in \tilde \Gamma
		 \end{cases}
\end{align*}
and show that it is continuous at $\bzero\in \tilde \Gamma$.
% Without loss of generality, we set $\bx = \bzero$. \Albert{[There is some clash of notation.  Later $\bx$ appears and definitely $\bx\neq\bzero$.]}
% \Albert{Remind the readers that our goal is to show \(\tilde\varphi\) is continuous at \(\bzero\) (which implies that \(f/g = \tilde\varphi \circ\Phi^{-1}\) is continuous at \(\ba\)).}
    

Since $\bQ$ is invertible on $\tilde \Gamma$, namely on a small region around $\bzero$, there is a strictly positive constant $c$ such that $|\bQ(\bx)\bx^\perp|\geq c |\bx^\perp|$ in some open neighborhood $U'$ with $\overline{U'}\subsetneq \tilde U$.
Then, we have for $\bx\in U'\setminus\tilde \Gamma$, 
    \begin{align*}
        \left|\tilde\varphi(\bx)-\tilde \varphi(\bzero)\right|
        & =\left|\frac{\tilde f(\bx)}{\tilde g(\bx)}- \lambda(\bzero)\right|\\
        &=\frac{| \bP(\bx)\bx^\perp -\Lambda(\bzero)\bQ(\bx)\bx^\perp|}{|\bQ(\bx)\bx^\perp|}\\
        & \leq  \frac{ \| \bP(\bx) -\Lambda(\bzero)\bQ(\bx)\|\cdot |\bx^\perp|}{c|\bx^\perp|}\\
        & = \frac{\| \bP(\bx) -\Lambda(\bzero)\bQ(\bx)\|}{c} \to 0 \text{ as }\bx\to \bzero.
    \end{align*}
    % \Albert{[In the 1st line \(\tilde\lambda\) is a complex number, but in the 2nd line \(\tilde\lambda\) is a 2x2 matrix. The equality needs to be checked carefully.]}
    Here $|\cdot|$ is the absolute value for a complex number and the standard Euclidean norm for an element of $\RR^2$, and  $\|\cdot\|$ is the operator norm defined by $\| A\|=\sup_{|\bz|=1} |A\bz|$ for a linear operator $A\coloneqq \RR^2 \to \RR^2$.  
   With the continuity of $\tilde\varphi|_{\tilde \Gamma}$,
   % we see that for any $\epsilon>0$, there is some $\delta>0$ such that $|\bx|<\delta$ implies $|\tilde\varphi(\bx)-\tilde \varphi(\bzero)|<\epsilon$ regardless if $\bx\in \tilde\Gamma$ or not.
   this shows the continuity of $\tilde\varphi$ at $\bzero$.  Hence \(f/g = \tilde\varphi \circ\Phi^{-1}\) is continuous at \(\ba\).
      % \Albert{[I can follow the proof.  But the implicit conversion between 2x2 matrices, 2D vector, and complex vector is a bit messy.]}

\vspace{5pt}
\noindent\emph{Proof of \hyperref[cl:P and Q]{Claim}.}  
We construct instances of $\bP,\bQ\in C^0(\tilde U, \RR^{2\times 2})$ with the stated properties.
For each $\bx\in\tilde U$, let us consider a function $F:[0,1]\to \CC\cong \RR^2$ by $F(t)= \tilde f(t\bx^\perp, \bx^\parallel)$. Since $F(0)=\tilde f(\bzero, \bx^\parallel)=0$, we have by the fundamental theorem of calculus that,
\begin{align*}
    \tilde f(\bx)=F(1)=\int^1_0 \frac{dF}{dt} dt 
    =\int^1_0 \nabla_{\bx^\perp}\tilde f(t\bx^\perp,\bx^\parallel) \cdot \bx^\perp dt. 
\end{align*}
 Using this relation, we define $\bP(\bx) = \int^1_0 \nabla_{\bx^\perp}\tilde f(t\bx^\perp,\bx^\parallel) dt$. Clearly $\tilde f(\bx)=\bP(\bx)\bx^\perp$. Note also that $\bP$ is continuous due to the $C^1$ regularity of $\tilde f$, and the invertibility of $\bP(\bzero,\bx^\parallel)=\nabla_{\bx^\perp}\tilde f(\bzero,\bx^\parallel)$ follows from the simple zero assumption of $f$. We can construct $\bQ$ using the same argument.


% For this we keep  using the same local coordinate system as above and show that $D\tilde f/D\tilde g(\bzero)$ is a complex number.  

% We claim that $D\tilde f(\bzero)=\tilde \varphi(\bzero) D\tilde g(\bzero)$ even when $\tilde \varphi=\Phi\circ \varphi$ is merely continuous but not necessarily differentiable at $\bzero$. Then it is clear from the definition that $D\tilde f/D\tilde g(\bzero)=\tilde \varphi(\bzero)$. This shows that $\tilde f$ and $\tilde g$ are complex linearly related, and hence so are $f$ and $g$.

% To verify the claim, recall that $D\tilde f(\bzero)$ is defined as a unique linear operator $L:\RR^n\to \CC$ such that 
% \begin{align*}\label{eq:def_of_Df}
%    \lim_{t\to 0} \frac{\tilde f(\bx(t))-\tilde f(\bzero) - L(\bx(t)) }{|\bx(t)|}=0
% \end{align*}
% for any path such that $\bx(0)=\bzero$.
%  Since $\tilde g$ is of $C^1$, we have $\tilde g(\bx)=Q(\bx)\bx^\perp$
%  % $\tilde g(\bx)=D\tilde g(\bzero)\bx + R_g(\bx)$ 
%  near $\bzero$
%  % with the remainder term $R_g\in o(|\bx|)$, 
%  and hence 
%  % $f(\bx(t))=(Dg(\bzero)\bx(t)+R_g(\bx))\varphi(\bx(t))$
%  $f(\bx(t))=(\bQ(\bzero)\bx^\perp(t))\varphi(\bx(t))$
%  . We now see that \eqref{eq:def_of_Df} is satisfied if and only if $L=\varphi(\bzero)Dg(\bzero)$.
%  \Sadashige{If this is too elementary and distractinge, we can simply say "It is easy to see that $Df(\bzero)=\varphi(\bzero) Dg(\bzero)$ when $\varphi$ is merely continuous but not necessarily differentiable at $\bzero$.".}
\end{proof}

% \section{First Order Taylor Expansion}
% \label{app:FirstOrderTaylorExpansion}
% The smoothness of l'H\^opital quotients of functions relies on results concerning the smoothness of the remainder term in Taylor's formula.  

% Suppose \(f\colon I\to\RR\) is a continuously differentiable function of one real variable over an open interval \(I\ni a\).  Then \(f\) can be expressed using the first-order Taylor formula as \(f(x) = f(a) + f'(a)(x-a) + (x-a)h(x)\), where the function \(h\) satisfies \(\lim_{x\to a}h(x) = 0\).  This form of the Taylor remainder is known as \emph{Peano's form}, first characterized by 
% Giuseppe Peano in 1889 \cite{Peano:1889:UNF,Persson:2017:HST}.   

% If we further assume that \(f\in C^\infty(I)\), then the Peano remainder function \(h\) is also \(C^\infty\)-smooth.  This regularity result is known as Hadamard's Lemma, which states that \(f(x) = f(a) + (x-a)g(x)\) for some \(g\in C^\infty(I)\).  An elementary proof of Hadamard's lemma can be found in \cite{Nestruev:2003:SMO}.

% More generally, if \(f\in C^k(I)\) for some \(k\geq 1\), then the first-order Peano remainder function \(h(x)\) (or \(g(x) = h(x)+f'(a)\) in the form of Hadamard's lemma) belongs to \(C^{k-1}(I)\).  This was shown by Hassler Whitney in 1943 \cite{Whitney:1943:DRT}.

% \begin{lemma}[Whitney 1943 \cite{Whitney:1943:DRT}]
% \label{lemma:Whitney}
%     Suppose \(f\in C^k(I)\) with \(k\geq 1\) on an open interval \(I\ni 0\).  Then there exists a function \(g\in C^{k-1}(I)\) such that
%     \(
%         f(x) = f(0) + xg(x),
%     \)     
%     with \(g^{(j)}(0) = {1\over j+1}f^{(j+1)}(0)\) for \(0\leq j\leq k-1\).
% \end{lemma}
% \begin{proof}
%     For \(x\neq 0\), the finite difference quotient \(g = {f(x) - f(0)\over x}\) can be written as \(g(x) = {1\over x}\int_0^xf'(t)\, dt\).   The derivatives of \(g(x)\) admit the following integral formula:
%     \begin{align*}
%     \label{eq:WhitneyMagicFormula}
%         {d^j\over dx^j}g(x) = {1\over x^{j+1}}\int_0^x t^j f^{(j+1)}(t)\, dt,\quad j=0,\ldots,k-1.
%     \end{align*}
%     This identity can be verified by repeated integration by parts and mathematical induction; we omit the details here.  
%     Since \(f\in C^k(I)\), it follows that \(g\in C^k(I\setminus\{0\})\subset C^{k-1}(I\setminus\{0\})\).
%     It remains to show that \(g^{(j)}(x)\) has a limit as \(x\to 0\), for \(0\leq j\leq k-1\).
%     Define \(L_j \coloneqq {1\over j+1}f^{(j+1)}(0)\), which can also be expressed as \(L_j = {1\over x^{j+1}}\int_0^x t^jf^{(j+1)}(0)\, dt \) for any \(x\neq 0\).  Then we compute:
%     \begin{align*}
%         \left|g^{(j)}(x) - L\right| & = 
%         \left|{1\over x^{j+1}}\int_0^x t^j\left(f^{(j+1)}(t) - f^{(j+1)}(0)\right)\, dt\right|\\
%         &\leq \underbrace{{1\over x^{j+1}}\int_0^x t^j\, dt}_{{1\over j+1}} \max_{t\in [0,x]}\left|f^{(j+1)}(t) - f^{(j+1)}(0)\right|.
%     \end{align*}
%     Since \(f^{(j+1)}\) is continuous at \(0\), this expression tends to zero as \(x\to 0\).  Therefore, each \(g^{(j)}\) extends continuously to \(x=0\), with \(g^{(j)}(0) = L_j\).  Hence, \(g\in C^{k-1}(I)\).
% \end{proof}

% Whitney also extended the result of \Cref{lemma:Whitney} to multidimensions in \cite{Whitney:1943:DRT}.  Suppose \(f\colon\RR^n\to\RR\) is defined in a neighborhood of the origin.  Then by fixing each \(x_2,\ldots,x_n\) and applying the same argument as \Cref{lemma:Whitney} on the variable \(x_1\), we obtain
% \begin{align*}
% \label{eq:WhitneyTowardsMultidiemsion}
%         f(x_1,x_2,\ldots,x_n) = f(0,x_2,\ldots,x_n) + x_1 g_1(x_1,x_2,\ldots,x_n)
% \end{align*}
% where \(f(0,x_2,\ldots,x_n)\) is still \(C^k\) and \(g_1\in C^{k-1}\).  The continuity of \({\partial^{i_1+\cdots+i_n}\over\partial x_1^{i_1}\cdots \partial x_n^{i_n}}g_1\) for \(i_1+\cdots+i_n\leq{k-1}\) is the result of applying \Cref{lemma:Whitney} on the function \({\partial^{i_2+\cdots +i_n}\over\partial x_2^{i_2}\cdots\partial x_n^{i_n}}f\).

% Repeating the process of \eqref{eq:WhitneyTowardsMultidiemsion} yields
% \begin{align*}
% \label{eq:WhitneyMultidimensionScalar}
%         f(\bx) = f(\bzero) + x_1g_1(x_1,\ldots,x_n) + x_2g_2(x_2,\ldots,x_n)+\cdots+ x_n g_n(x_n)
% \end{align*}
% where \(g_i\in C^{k-1}\) and \(g_i(\bzero) = {\partial f\over\partial x_i}(\bzero)\).

% The result \eqref{eq:WhitneyMultidimensionScalar} also holds for vector-valued functions \(\bff = \begin{bsmallmatrix}
%     f_1\\|\\f_m
% \end{bsmallmatrix}\colon\RR^n\to\RR^m\) by applying \eqref{eq:WhitneyMultidimensionScalar} componentwise, which is concluded the following lemma.
% In the lemma we single out the factors \(\sum_{i}{\partial\bff\over\partial x_i}(\bzero)x_i = D\bff(\bzero)\bx\) to resemble Peano's form of Taylor expansion.

% \begin{lemma}[First-order Taylor formula with regularity]
% \label{lemma:Taylor}
% Let \(\bff\colon U\subset\RR^n\to\RR^m\) be a \(C^k\)-continuous vector-valued function on a neighborhood \(U\) of a given point $\ba\in \RR^n$.
% Then there exists \(\bH\in C^{k-1}(U;\RR^{m\times n})\) matrix-valued function so that \(\bH(\ba) = \bzero\) and
% \begin{align*}
%     \bff(\bx) = \bff(\ba) + D\bff(\ba)(\bx-\ba) + \bH(\bx)(\bx-\ba),\quad\bx\in U.
% \end{align*}
% \end{lemma}

% \section{Proof of Theorem \ref{thm:RealQuotient1D}}
% \label{app:ProofOfRealQuotient1D}



% % \begin{lemma}
% %     Let \(f\colon I\to\FF\), with \(\FF = \RR\) or \(\CC\), be a \(C^k\)-smooth function of one variable over a real open interval \(I\), with \(k\geq 1\).  Suppose \(f(a) = 0\) at \(a\in I\) and \(f'(a) \neq 0\). Then there exists a neighborhood \(U\subset I\) of \(a\) where \(f(x) = f'(a)(x-a)F(x)\) for some \(C^{k-1}\)-smooth function \(F\colon U\to\FF\) satisfying \(F(a) = 1\) and \(F(x)\) is non-vanishing on \(U\).
% % \end{lemma}
% % \begin{proof}
% %     By the Taylor Theorem with Peano's remainder, there exists a function \(h(x)\in C^{k-1}(I;\FF)\) such that \(f(x) = f'(a)(x-a) + (x-a)h(x)\) with \(\lim_{x\to a}h(x) = 0\).  Note that the classical Taylor's Theorem which applies to real-valued function extends to complex-valued function by applying Taylor's Theorem on real and imaginary part separately.  Taking the common factor \(f'(a)(x-a)\) out to obtain \(f(x) = f'(a)(x-a)(1+h(x)/f'(a))\) and define \(F(x) = (1+h(x)/f'(a))\in C^{k-1}(I;\FF)\).  Therefore \(f(x) = f'(a)(x-a)F(x)\).  The limit of \(F(x)\) at \(x\to a\) is 1.  By continuity of \(F\) we have \(F(a)=1\) and there exists a nieghborhood of \(a\) in which \(F\) is nonzero.
% % \end{proof}
% \begin{proof}[Proof of \autoref{thm:RealQuotient1D}]
%     Let \(\Gamma\) be the zero set of \(f\in C^k(I)\) and \(g\in C^k(I)\).  For \(x\in I\setminus\Gamma\) we define \(\varphi\) by \(\varphi(x) = f(x)/g(x)\), which belongs to \(C^{k}(I\setminus\Gamma)\subset C^{k-1}(I\setminus\Gamma)\) and non-vanishing.  It remains to show that \(\varphi\) is \(C^{k-1}\)-smoothly defined in a neighborhood of \(\Gamma\).  
%     By \Cref{lemma:Whitney}, near a simple zero \(a \in \Gamma\) of both \(f\) and \(g\), we can express  
% \[
% f(x) = f'(a)(x-a)F(x), \quad g(x) = g'(a)(x-a)G(x)
% \]
% for \(x\) in a neighborhood \(U\) of \(a\), where \(F,G\in C^{k-1}(U)\), \(F(a) = G(a) = 1\), and both \(F(x)\) and \(G(x)\) are nonzero on \(U\). 
% Thus, for \(x \in U\), we define  
% \[
% \varphi(x) \coloneqq \frac{f'(a)}{g'(a)} \cdot \frac{F(x)}{G(x)},
% \]
% which is smooth, non-vanishing,  $\varphi(a)=f'(a)/g'(a)$, and agrees with \(f(x)/g(x)\) on \(U\setminus\{a\}\).  
% \end{proof}

% \section{Proof of Theorem \ref{thm:RealQuotientnD}}
% \label{app:ProofOfRealQuotientnD}
% \begin{proof}[Proof of \autoref{thm:RealQuotientnD}]
%     Let \(\Sigma\) be the common zero set of \(f,g\in C^k(\Omega)\).  On \(\Omega\setminus\Sigma\) define \(\bx\in\Omega\setminus\Sigma\) define \(\varphi(\bx) = f(\bx)/g(\bx)\), which belongs to \(C^{k}(\Omega\setminus\Sigma)\) and is non-vanishing.
%     It remains to show that \(\varphi\) is \(C^{k-1}\)-continuously defined in a neighborhood of any point of \(\Sigma\).  Let \(\ba\in\Sigma\).  By the implicit function theorem, there exists a neighborhood \(U\) of \(\ba\), a neighborhood \(\tilde U\subset\RR^n\) of the origin \(\bzero\in\RR^n\), and a \(C^k\) bijection \(\Phi\colon \tilde U\xrightarrow{\simeq} U\) such that \(\Phi(\bzero) = \ba\) and \(\tilde\Sigma = \Phi^{-1}(\Sigma)\) is the coordinate hyperplane \(\tilde\Sigma = \{(x_1,\ldots, x_n)\,|\,x_1 = 0\}\cap\tilde U\).
%     Define \(\tilde f = f\circ\Phi\) and \(\tilde g = g\circ\Phi\), both of which are \(C^k\) functions on \(\tilde U\) and have simple zeros on \(\tilde\Sigma\).  
%     In this coordinate system, \(\nabla \tilde f(0,x_2,\ldots,x_n) = {\partial \tilde f\over\partial x_1}(0,x_2,\ldots,x_n)\be_1\) and \(\nabla \tilde g(0,x_2,\ldots,x_n) = {\partial \tilde g\over\partial x_1}(0,x_2,\ldots,x_n)\be_1\), where \(\be_1\) is the coordinate vector \((1,0,\ldots,0)\), and \({\partial\tilde f\over\partial x_1}\) and \({\partial\tilde g\over\partial x_1}\) are nonvanishing \(C^{k-1}\) function on \(\tilde\Sigma\).
%     By \eqref{eq:WhitneyTowardsMultidiemsion}, there exists \(p,q\in C^{k-1}(\tilde U)\) such that for \((x_1,\ldots,x_n)\in \tilde U\)
%     \begin{align*}
%         \tilde f(x_1,\ldots,x_n) = x_1 p(x_1,\ldots,x_n),\quad\tilde g(x_1,\ldots,x_n) = x_1q(x_1,\ldots,x_n)
%     \end{align*}
%     and \(p = {\partial\tilde f\over\partial x_1}\neq 0\) and \(q = {\partial\tilde g\over\partial x_1}\neq 0\) on \(\tilde\Sigma\).
%     Therefore there exists a neighborhood of \(\tilde\Sigma\), without loss of generality \(\tilde U\), on which \(p\) and \(q\) are nonvanishing.  Therefore, \(\tilde f\over\tilde g = {p\over q}\) is nonvanishing \(C^{k-1}\) continuous function on \(\tilde U\) and agrees with \(\varphi\circ \Phi\) on \(\tilde U\).  Thus \(\varphi\) extends to a \(C^{k-1}\)-continuous function in a neighborhood of \(\ba\).
% \end{proof}

% \section{Proof of Theorem \ref{thm:PathDependentlHopital}}\label{app:proof_PathDependentlHopital}

% To prove the theorem, we use a complex version of Theorem \ref{thm:RealQuotient1D}, which can be obtained with the complex  version of Lemma 1. For completeness, we state them as lemmas.

% \begin{lemma}
% \label{lemma:ComplexWhitney}
%     Suppose \(f\in C^k(I,\CC)\) with \(k\geq 1\) on an open interval \(I\ni 0\).  Then there exists a function \(g\in C^{k-1}(I)\) such that
%     \(
%         f(x) = f(0) + xg(x),
%     \)     
%     with \(g^{(j)}(0) = {1\over j+1}f^{(j+1)}(0)\) for \(0\leq j\leq k-1\).
% \end{lemma}
% \begin{proof}
%     We simply apply Lemma \ref{lemma:Whitney} to the real and the imaginary part of $f$ to get two functions, which we can use as the real and the imaginary part of $g$.
% \end{proof}

% \begin{lemma}[Smooth quotient of complex functions on a real line]\label{lem:ComplexQuotient1D}
%      Let \(f,g\in C^\infty(I,\CC)\) be smooth complex-valued functions on a real interval \(I\). Suppose their zeros are identical and their derivatives $f'$ and $g'$ are non-zero on these zero sets. Then there exists a non-vanishing smooth complex function \(\varphi\in C^\infty(I)\) such that \(f = g\varphi\).
% \end{lemma}
% \begin{proof}
%     The proof is identical with the proof for Theorem \ref{thm:RealQuotient1D}. 
% \end{proof}

% \begin{proof}[Proof of Theorem \ref{thm:PathDependentlHopital}]
%     Let us consider smooth functions $f_\gamma(t)\coloneqq f(\gamma(t))$ and $f_\gamma(t)=g(\gamma(t))$ on the interval $(-\epsilon,\epsilon)$. The quotient $\varphi\coloneqq  f_\gamma/ g_\gamma$ is smooth and nowhere vanishing with $\varphi(0)=f'_\gamma(0)/g'_\gamma(0)$ due to Lemma \ref{lem:ComplexQuotient1D}. Finally noting that $g_\gamma'(0)=Dg|_{\gamma(0)}\gamma'(0)$ and $f_\gamma'(0)=Df|_{\gamma(0)}\gamma'(0)= \frac{Df}{Dg} g_\gamma'(0)$, we obtain the stated expression.
% \end{proof}

% \section{Proof of Theorem \ref{thm:ComplexQuotientnD}}\label{app:proof_ ComplexQuotientnD}

% % \begin{proof}[Proof of \autoref{thm:ComplexQuotientnD}]
% %     We first show that $(ii)$ implies $(i)$. The continuity of $\varphi=f/g$ on $\Omega\setminus f^{-1}(0)$ is trivial. We prove now $\varphi$ is defined and continuous on at each point $p$ of $f^{-1}(0)$. In a small open neighborhood $B$ of $p$ we take a local coordinate $(\bx, \by)=(x_1,x_2, y_1,\ldots, y_{n-2})$ where $f^{-1}(0)\subset B=\{(\bx,\by)\in B \,\mid\, \bx=\mathbf{0}\}$. 

% %     Note that for each $\by$, consider the slice $B_\by$ of $B$. On $B_\by$ 

% % \TODO{To be continued.}

% % And $\varphi(\bzero,\by)$ is continuous in $\by$ since $\varphi(\bzero,\by)=Df/Dg(\bzero,\by)\in \CC$ for each $\by$ and the matrix-valued function $Df/Dg$ is smooth everywhere by assumption.

% % Now $\varphi$ may appear to be continuous at $(\bzero,\by)$ as $\varphi(\bzero,\by)$ is continuous in $\by$, and for each curve $\gamma$ transversally intersecting with $\{\bx=0\}$ at $\gamma(0)=(\bzero,\by)$, $\varphi(\gamma(s))$ is continuous and $\varphi(\gamma(0))=\varphi(\bzero,\by)$. This however does not exclude the possibility of the existence of a path $\gamma$ that tangentially touches $\{\bx=0\}$ at $(\bzero,\by)=\gamma(0)$ and $\lim_{s\to 0}\varphi(\gamma(s))\neq \varphi(\bzero,\by)$.


% % \TODO{To be continued.}

% %     We want to use the triangle inequality
% %     \begin{align*}
% %         |\varphi(\bx, \by)-\varphi(\bzero,\bzero)|\leq |\varphi(\bzero,\by)-\varphi(\bzero,\bzero)|+|\varphi(\bzero,\by)-\varphi(\bx,\by)|.
% %     \end{align*}
% %     The control of the first term in the right-handed-side is easy since $\varphi(\bzero,\by)$ is continuous in $\by$. However, controlling the second term is non-trivial. For each given $\epsilon>0$, there exists $\delta_\by>0$ for which $|\bx|<\delta_\by$ implies $|\varphi(\bzero,\by)-\varphi(\bx,\by)|<\epsilon$. However, the strict positivity $\inf_{\by\in I}\delta_\by >0$ on some small domain $I=[-\by_0,\by_0]^{n-2}$ is unclear at this point. Namely, we have to show the equicontinuity of $M\coloneqq\{\varphi_\by\}_{\by\in I}$. 

% %     To make our setting applicable for the (arguably) converse direction of the Ascoli-Arzela theorem,
% %     we consider the restriction of each $\varphi_\by$ on a small closed box $X=[-x_0,x_0]^2$ so that they are bounded and belong to Banach space
% %     $(C(X),\|\cdot\|_{\sup})$. For the restricted functions, we still use the same notation $\varphi_\by$ for simplicity. 

% %     We now show that $M$ is compact in $C(X)$. For any sequence $\{\varphi_{n}\}_n$ in $M$ is of the form $\{\varphi_{\by_n}\}_n$ and there is a subsequence $\by_{n_j}$ converging to some $y\in I$ as $I$ is compact. Then $\lim_{j->\infty}\|\varphi_{\by_{n_j}}\to \varphi_\by\|_{\sup}$ as  $\varphi_{\by_{n_j}}(\bx)\to \varphi_\by(\bx)$ for $\bx\neq \bzero$ due to the continuity of $\varphi$ on $\{\by \neq \bzero\}$ and for $\bx=\bzero$ due to the continuity of $\varphi(\bzero,\by)$ in $\by$.
    
% %     there is a convergent subsequence $f_{y_{n_j}}\to f_y$ for some $y$ 
    
% % \end{proof}

% % \Sadashige{I commented out my long proof below.}

% % The Peano reminder $R(\bx)\coloneqq H(\bx)\cdot \bx$ in Lemma 2 is in $o(|\bx|)$. That is, 
% % \begin{align*}
% %     \lim_{\bx\to \bzero} \frac{R(\bx)}{|\bx|}=0
% % \end{align*}
% % where $| \bx|=\sqrt{\sum_i x_i^2}$ is the standard Euclidean norm.

% % In our case, we have a stronger result. Let us assume the setting of Theorem 4. Then around each point of the zero set $\Gamma$, we may take without loss of generality a local coordinate system $\bx=(\bx^\perp, \bx^\parallel)$ with $2$ dimensional $\bx^\perp$ and $n-2$ dimensional $\bx^\parallel$ so that  $\Gamma$ is identical with the local $n-2$ dimensional hyperplane $\{\bx=(\bx^\perp,\bx^\parallel) \mid\, \bx^\perp=\bzero\}$. Here and in the rest of the appendix, the dimension of the zero vector $\bzero$ depends on the context. Then, we have the following:
% % \begin{lemma}[Refinement of Lemma \ref{lemma:Taylor}]
% % In the above local coordinates, we have $R\in o(|\bx^\perp|)$, that is,
% % \begin{align*}
% %     \lim_{\bx\to \bzero} \frac{R(\bx)}{|\bx^\perp|}=0.
% % \end{align*}
% % \end{lemma}

% % \begin{proof}
% % In our local coordinates, we have
% % \[
% % Df(\bx)=
% % \begin{pmatrix}
% % \begin{array}{ccccc}
% % \multicolumn{2}{c}{\multirow{2}{*}{\scalebox{1.2}{$DF$}}}
% %   & \multicolumn{3}{c}{\multirow{2}{*}{\scalebox{1.2}{$\ RF\ $}}} \\
% %  &  &  &  &  \\
% % \end{array}
% % \end{pmatrix}
% % \begin{pmatrix}
% % \begin{array}{cc}
% % \bx^\perp\\
% % \bx^\parallel
% % \end{array}
% % \end{pmatrix}
% % \]
% % where $DF$ is a $2$ by $2$ matrix valued function, which is invertible on the simple zeros $\Gamma$, and $RF$ is a $2$ by $(n-2)$ matrix valued function, which vanishes on $\Gamma$.

% % Lemma \ref{lemma:Taylor} states
% % \begin{align*}
% %    f(\bx)=f(\bzero)+Df(\bzero)\cdot \bx + H(\bx)\cdot \bx
% % \end{align*}
% % where $H(\bx)\to \bzero$ as $\bx\to \bzero$. Hence we have for $R(\bx)=H(\bx)\cdot \bx$ that, 
% % \begin{align*}
% %     R(\bx) &= f(\bx)- f(\bzero)- Df(\bzero)\cdot (\bx)\\
% %     &= f(\bx^\perp,\bx^\parallel) - f(\bzero,\bx^\parallel) - DF(\bzero)\cdot \bx^\perp\\
% %     &=\left( DF(t_{\bx} \bx^\perp,\bx^\parallel) - DF(\bzero) \right)\cdot \bx^\perp
% % \end{align*}
% % where $t_{\bx}\in [0,1]$. Here we used $f(\bzero)=f(\bzero,\bx^\parallel)=0$ and applied the mean value theorem to the function $\tau:[0,1]\to \CC$ given by $\tau(t)\coloneqq f(t \bx^\perp,\bx^\parallel)$. 

% % Since $DF(t_{\bx}\bx^\perp,\bx^\parallel)\to DF(\bzero)$ as $\bx\to \bzero$, we have
% % \begin{align*}
% %     \lim_{\bx\to \bzero}\frac{R(\bx)}{|\bx^\perp|}
% %     \leq  \lim_{\bx\to \bzero}\frac{\| DF(t_\bx \bx^\perp,\bx^\parallel) - DF(\bzero,\bx^\parallel) \|\,|\bx^\perp|}{|\bx^\perp|}\to 0
% % \end{align*}
% % where $\|\cdot \|$ is the operator norm defined by $\|A\|=\sup_{|x|=1}|Ax|$ for a linear operator $A:\RR^2\to \CC\simeq \RR^2$. 
% % \end{proof}

% In a 2D domain, \autoref{thm:ComplexQuotientnD} is a direct consequence of \autoref{thm:PathDependentlHopital}. 
% If \(f\) and \(g\) are complex linear, 
% with  the quotient of derivative
% % in \eqref{eq:PathDependentlHopital} is a scaled rotation 
% \(({Df\over Dg}) = \begin{bsmallmatrix}
%     u&-v\\
%     v & u
% \end{bsmallmatrix}\) with some $u,v\in \RR$ at each zero, 
% then \eqref{eq:PathDependentlHopital} simplifies to 
% \begin{align*}
%     \lim_{t=0}{f(\gamma(t))\over g(\gamma(t))} = {1\over g'}(u+\ii v)g' = u+\ii v
% \end{align*}
% which is a single complex number that depends only on the point \(\gamma(0)\) and is independent of the path direction \(\gamma'\). This gives a proof of \autoref{thm:ComplexQuotientnD} in 2D. 

% In higher dimensions, however, \autoref{thm:ComplexQuotientnD} does not automatically follow. From the considerations above, we see that when $f$ and $g$ are smooth and complex linear, the quotient $\varphi$  is continuous along any path intersecting $\Gamma$ transversally. In addition, $\varphi$ is smooth when restricted to  $\varphi|_{\Gamma}$ or to $\varphi_{\Omega\setminus\Gamma}$. These properties are, however, not enough to conclude the continuity of $\varphi$. In fact, we can construct a function that satisfies all these properties that is yet undefined on $\Gamma$, as in the following example. 
% To prove the theorem, we need to rely on the structures specific to the quotient of $f$ and $g$.\newline
% % 
% \vspace{-10pt}
% \begingroup
% % \begin{wrapfigure}{r}{0.15\textwidth}
% % \centering
% % % trim={left bottom right top}
% % \includegraphics[trim={2cm 0 2cm 0},clip, width=0.12\textwidth]{images/bump_functs.png}
% % \caption{test.}
% % \label{fig:bump_functions}
% % \end{wrapfigure}
% % Adjusted settings
% \setlength{\intextsep}{10pt}%
% \setlength{\columnsep}{10pt}%
% % \vspace{3cm}
% \begin{wrapfigure}{r}{0.2\textwidth}
% % % trim={left bottom right top}
%   % \centering\includegraphics[width=1.0\linewidth,trim={13cm 1cm 13cm 0}, clip]{images/bump_functions_no_labels.png}
%   % \caption{Basic layout}
%       \begin{picture}(0,150)
%         \put(0,0){\includegraphics[width=.2\columnwidth,trim={13cm 1cm 13cm 0}, clip]{images/bump_functions_no_labels.png}}
%         \put(45,22){\small\(x\)}
%         \put(16,80){\small\(z\)}
%     \end{picture}
%     % \caption{ Cut-off view of $xz$-plane. Discs are the supports of the bump functions $b_n$, lying on the curve $z=2\sqrt{x}$. }
% \end{wrapfigure}
% \begin{example}\label{eg:infinite_bump_functions}
%     Consider a collection of smooth bump functions $\{b_n\}_{n\in \NN}$ in $\RR^3$. Here each $b_n$ takes its maximum value $1$ at the center $(2^{-n},0, 2^{1-{n\over 2}})$ of its support given as an open ball of radius $2^{-n}$. Their centers lie on the curve $z=2\sqrt{x}$, and the supports intersect $z$-axis, as visualized in the right image.
%     Using these bump functions, let us define $\varphi(x,y,z)= \sum_{n=0}^\infty b_n(x,y,z)$. 
%     Notice that  any straight line on the $xz$-plane passing through the origin intersects only finitely many of these open balls. Hence, along such lines $\gamma:(-\epsilon,\epsilon)\to \RR^3$ with $\gamma(0) = \mathbf{0}$, we have $\lim_{t\to0} \varphi(\gamma(t))=0$. However, any curve $(t,0,a \sqrt{t})$ with $a>2$, intersects these balls infinitely many times as $t\to 0$, making $\varphi$ discontinuous at the origin.
% \end{example}
%  \endgroup
 
% We now prove \autoref{thm:ComplexQuotientnD}. The key component is the simple zero condition of $f$ and $g$. This ensures that the leading order of these functions on $\Omega\setminus \Gamma$ does not involve the component of the local coordinates tangent to $\Gamma$. This leads to the continuity of the quotient along any path approaching to $\Gamma$, possibly non-transversalilly, as in Example \ref{eg:infinite_bump_functions}.

% \begin{proof}[Proof of \autoref{thm:ComplexQuotientnD}]
% {$(i)\to (ii)$.} We show that $Df/Dg(\ba)$ is a complex number at any $\ba\in \Gamma$.
% It follows from the $C^1$ regularity of $g$ and the continuity of $\varphi$ that $Df(\ba)=\varphi(\ba)Dg(\ba)$ even if $\varphi$ is not differentiable. 
% % 
% \suggest{
% To see this, let us recall that $D f(\ba)$ is defined as a unique linear operator $L:\RR^n\to \CC$ such that 
% \begin{align*}\label{eq:def_of_Df}
%    \lim_{\bx \to \ba} \frac{f(\bx)-f(\ba) - L(\bx-\ba) }{|\bx-\ba|}=0.
% \end{align*}
% Since $g(\bx)=Dg(\ba)(\bx-\ba)+H(\bx)(\bx-\ba)$ for some continuous function $H$ such that $H(\ba)=0$ due to Lemma \ref{lemma:Taylor}, we have
% \begin{align*}\label{eq:def_of_Df}
%    \frac{f(\bx)-f(\ba) - L(\bx) }{|\bx|}
%    &= \frac{|\varphi(\bx)(Dg(\ba)+H(\bx-\ba))(\bx-\ba)-L(\bx-\ba)|}{|\bx-\ba|}\\
%    &\leq \|\varphi(x)(Dg(\ba)+H(\bx-\ba))-L\|
% \end{align*}
% where $\|\cdot\|$ denotes the operator norm. This converges to $0$ as $\bx\to \ba$ if $L=\varphi(\ba)Dg(\ba)$.
% }\Sadashige{I think we may drop the above part as it is almost trivial.}

% Then we obtain the claimed statement 
% % $Df/Dg(\ba)=\varphi(\ba)\in \CC$
% by the definition of $Df/Dg$.

% {$(ii)\to (i)$.}
%     Let \(\Gamma\) be the common zero set of \(f,g\in C^1(\Omega;\CC)\).  
%     The quotient \(\varphi(\bx) = f(\bx)/g(\bx)\) for \(\bx\in \Omega\setminus\Gamma\), is \(C^1\)-continuous and nonvanishing.
%     It remains to show that \(f/g\) extends to a \(C^{0}\)-continuous function in a neighborhood of any point \(\ba\in \Gamma\).  

%     % As in the proof of \autoref{thm:RealQuotientnD}, we may take without loss of generality a local coordinate system centered at $\ba$ and $\Gamma$ is 
%     % we may take without loss of generality a local coordinate system $\bx=(\bx^\perp, \bx^\parallel)$ centered at $\ba$ with $2$ dimensional $\bx^\perp$ and $n-2$ dimensional $\bx^\parallel$ so that  $\Gamma$ is identical with the local $n-2$ dimensional hyperplane $\{\bx=(\bx^\perp,\bx^\parallel) \mid\, \bx^\perp=\bzero\}$.

    
%     As in the proof of \autoref{thm:RealQuotientnD}, the implicit function theorem assures the existence of a neighborhood \(U\) of \(\ba\), a neighborhood \(\tilde U\) of the origin \(\bzero\in\RR^n\), and a \(C^1\) bijection \(\Phi\colon\tilde U\xrightarrow{\simeq} U\) such that \(\Phi(\bzero) = \ba\) and \(\tilde\Gamma = \Phi^{-1}(\Gamma)\) is the subspace \(\tilde\Gamma = \{\bx=(\bx^\perp, \bx^\parallel)\,|\,  \bx^\perp=\bzero\}\cap\tilde U\). Here and in the rest of the proof, we write $\bx^\perp=(x_1,x_2)$ and $\bx^\parallel=(x_3,\ldots,x_n)$ for simplicity.
    
%     Now let \(\tilde f = f\circ\Phi\) and \(\tilde g = g\circ\Phi\), which are \(C^1\) functions vanishing on \(\tilde\Gamma\).
%     Then there exists invertible \(\bP,\bQ\in C^{0}(\tilde U;\RR^{2\times 2})\) such that 
%     \begin{align*}
%         \tilde f(\bx) = \bP(\bx)\bx^\perp,\quad
%         \tilde g(\bx) = 
%         \bQ(\bx)\bx^\perp,\quad \bx^\perp\coloneqq (x_1,x_2)
%     \end{align*}
%     with $\bP(\bzero,\bx^\parallel)=\nabla_{\bx^\perp}\tilde f(\bzero,\bx^\parallel)$ and $\bQ(\bzero,\bx^\parallel)=\nabla_{\bx^\perp}\tilde g(\bzero,\bx^\parallel)$ being invertible since $(\bzero,\bx^\parallel)$ are simple zeros.

% It follows from the continuity of $\bP,\bQ$ and the assumption that \(f\) and \(g\) are complex linearly related that there exists a continuous complex-valued function \(\tilde \lambda\) on $\tilde \Gamma$ so that
% % \(Df = \lambda Dg\) and
% % \(D\tilde f = \tilde\lambda D\tilde g\) where \(\tilde \lambda = \lambda\circ\Phi\). 
%      \(\bP = \tilde\lambda\bQ\) on \(\tilde\Gamma\).  
% % Note that $\tilde \lambda= \bP \bQ^{-1}
% % =\nabla_{\bx^\perp}\tilde f(\bzero,\bx^\parallel) \nabla_{\bx^\perp}\tilde g(\bzero,\bx^\parallel)^{-1}
% % $ is continuous since $\bP, \bQ$ are continuous on $\tilde \Gamma$.
%     % 
% Using this, we define a function $\tilde \varphi$ on $\tilde \Gamma$ by
%     \begin{align*}
%     \tilde\varphi(\bx)=
%     \begin{cases}
% 			\tilde f(\bx)/\tilde g(\bx), & \bx \notin \tilde\Gamma,\\
%             \tilde \lambda(\bx^\parallel), &  \bx \in \tilde \Gamma
% 		 \end{cases}
% \end{align*}
% and show that it is continuous at any $\bx\in \tilde \Gamma$. Without loss of generality, we set $\bx = \bzero$.
    

% Since $\bQ$ is invertible in $\tilde U$, there is a strictly positive constant $c$ such that $|\bQ(\bx)\bx^\perp|\geq c |\bx^\perp|$ in some open neighborhood $U'$ with $\overline{U'}\subsetneq \tilde U$.
% Then, we have for $\bx\in U'\setminus\tilde \Gamma$, 
%     \begin{align*}
%         \left|\tilde\varphi(\bx)-\tilde \varphi(\bzero)\right|
%         & =\left|\frac{\tilde f(\bx)}{\tilde g(\bx)}-\tilde \lambda(\bzero)\right|\\
%         &=\left| \frac{ \bP(\bx)\bx^\perp -\tilde\lambda(\bzero)\bQ(\bx)\bx^\perp}{\bQ(\bx)\bx^\perp}\right|\\
%         & \leq  \frac{ \| \bP(\bx) -\tilde\lambda(\bzero)\bQ(\bx)\|\cdot |\bx^\perp|}{c|\bx^\perp|}\\
%         & = \frac{\| \bP(\bx) -\tilde\lambda(\bzero)\bQ(\bx)\|}{c} \to 0 \text{ as }\bx\to \bzero
%     \end{align*}
%     where $\|\cdot\|$ is the operator norm defined by $\| A\|=\sup_{|\bz|=1} |A\bz|$ for a linear operator $A\coloneqq \RR^2 \to \RR^2$.  
%    With the continuity of $\tilde\varphi$ on $\tilde \Gamma$,
%    % we see that for any $\epsilon>0$, there is some $\delta>0$ such that $|\bx|<\delta$ implies $|\tilde\varphi(\bx)-\tilde \varphi(\bzero)|<\epsilon$ regardless if $\bx\in \tilde\Gamma$ or not.
%    this shows the continuity of $\tilde\varphi$, and hence of $\varphi$ at $\bzero$.

% % For this we keep  using the same local coordinate system as above and show that $D\tilde f/D\tilde g(\bzero)$ is a complex number.  

% % We claim that $D\tilde f(\bzero)=\tilde \varphi(\bzero) D\tilde g(\bzero)$ even when $\tilde \varphi=\Phi\circ \varphi$ is merely continuous but not necessarily differentiable at $\bzero$. Then it is clear from the definition that $D\tilde f/D\tilde g(\bzero)=\tilde \varphi(\bzero)$. This shows that $\tilde f$ and $\tilde g$ are complex linearly related, and hence so are $f$ and $g$.

% % To verify the claim, recall that $D\tilde f(\bzero)$ is defined as a unique linear operator $L:\RR^n\to \CC$ such that 
% % \begin{align*}\label{eq:def_of_Df}
% %    \lim_{t\to 0} \frac{\tilde f(\bx(t))-\tilde f(\bzero) - L(\bx(t)) }{|\bx(t)|}=0
% % \end{align*}
% % for any path such that $\bx(0)=\bzero$.
% %  Since $\tilde g$ is of $C^1$, we have $\tilde g(\bx)=Q(\bx)\bx^\perp$
% %  % $\tilde g(\bx)=D\tilde g(\bzero)\bx + R_g(\bx)$ 
% %  near $\bzero$
% %  % with the remainder term $R_g\in o(|\bx|)$, 
% %  and hence 
% %  % $f(\bx(t))=(Dg(\bzero)\bx(t)+R_g(\bx))\varphi(\bx(t))$
% %  $f(\bx(t))=(\bQ(\bzero)\bx^\perp(t))\varphi(\bx(t))$
% %  . We now see that \eqref{eq:def_of_Df} is satisfied if and only if $L=\varphi(\bzero)Dg(\bzero)$.
% %  \Sadashige{If this is too elementary and distractinge, we can simply say "It is easy to see that $Df(\bzero)=\varphi(\bzero) Dg(\bzero)$ when $\varphi$ is merely continuous but not necessarily differentiable at $\bzero$.".}
% \end{proof}
\end{document}
